\documentclass[12pt,a4paper]{report}

% ─── Packages ───────────────────────────────────────────────────────────────
\usepackage[utf8]{inputenc}
\usepackage[T1]{fontenc}
\usepackage[french]{babel}
\usepackage{geometry}
\usepackage{xcolor}
\usepackage{listings}
\usepackage{fancyhdr}
\usepackage{titlesec}
\usepackage{hyperref}
\usepackage{enumitem}
\usepackage{mdframed}
\usepackage{booktabs}
\usepackage{array}
\usepackage{parskip}
\usepackage{graphicx}
\usepackage{tcolorbox}
\usepackage{fontawesome5}
\usepackage{tikz}
\usepackage{amsmath}
\usepackage{microtype}
\usepackage{etoolbox}
\usepackage{caption}
\usepackage{multicol}
\usepackage{tabularx}
\usepackage{colortbl}
\usepackage{lmodern}
\usepackage{setspace}
\usepackage{ifthen}
% CORRECTION : zref-abspage pour éviter les conflits de compteurs avec cleardoublepage
\usepackage{emptypage}

\tcbuselibrary{skins, breakable, listings}

% ─── Mise en page ───────────────────────────────────────────────────────────
\geometry{
  top=2.5cm, bottom=2.8cm,
  left=2.8cm, right=2.5cm,
  headheight=18pt
}
\setlength{\headsep}{10pt}

% ─── Palette de couleurs ────────────────────────────────────────────────────
\definecolor{bleuENSA}{RGB}{0, 56, 101}
\definecolor{bleuClair}{RGB}{214, 228, 243}
\definecolor{bleuMoyen}{RGB}{70, 130, 180}
\definecolor{bleuAccent}{RGB}{0, 102, 153}
\definecolor{bleuTitre}{RGB}{0, 78, 140}
\definecolor{vertCode}{RGB}{0, 100, 0}
\definecolor{vertSucces}{RGB}{0, 110, 50}
\definecolor{vertClairBg}{RGB}{230, 247, 237}
\definecolor{grisCode}{RGB}{248, 249, 250}
\definecolor{grisClairBg}{RGB}{245, 246, 248}
\definecolor{grisFond}{RGB}{250, 251, 253}
\definecolor{orangeNote}{RGB}{180, 90, 0}
\definecolor{orangeClairBg}{RGB}{255, 243, 215}
\definecolor{rougeErr}{RGB}{140, 15, 15}
\definecolor{rougeClairBg}{RGB}{253, 232, 232}
\definecolor{grisBorder}{RGB}{180, 195, 210}
\definecolor{grisTexte}{RGB}{70, 80, 90}
\definecolor{violetInfo}{RGB}{80, 45, 140}
\definecolor{violetClairBg}{RGB}{238, 230, 255}
\definecolor{orDore}{RGB}{153, 117, 0}
\definecolor{sujBg}{RGB}{240, 248, 255}
\definecolor{sujBorder}{RGB}{0, 102, 153}

% ─── Style de code Java ─────────────────────────────────────────────────────
\lstdefinestyle{java}{
  language=Java,
  basicstyle=\ttfamily\footnotesize,
  keywordstyle=\color{bleuENSA}\bfseries,
  stringstyle=\color{rougeErr},
  commentstyle=\color{vertCode}\itshape,
  numberstyle=\tiny\color{gray},
  numbers=left,
  numbersep=10pt,
  stepnumber=1,
  backgroundcolor=\color{grisCode},
  frame=single,
  framerule=0.4pt,
  rulecolor=\color{grisBorder},
  breaklines=true,
  breakatwhitespace=false,
  tabsize=4,
  showstringspaces=false,
  captionpos=b,
  xleftmargin=18pt,
  xrightmargin=5pt,
  aboveskip=8pt,
  belowskip=8pt,
  extendedchars=true,
  inputencoding=utf8,
  literate=
    {é}{{\'e}}1 {è}{{\`e}}1 {ê}{{\^e}}1 {ë}{{\"e}}1
    {à}{{\`a}}1 {â}{{\^a}}1 {ä}{{\"a}}1
    {î}{{\^i}}1 {ï}{{\"i}}1
    {ô}{{\^o}}1 {ö}{{\"o}}1
    {ù}{{\`u}}1 {û}{{\^u}}1 {ü}{{\"u}}1
    {ç}{{\c{c}}}1
    {É}{{\'E}}1 {È}{{\`E}}1 {Ê}{{\^E}}1
    {À}{{\`A}}1 {Â}{{\^A}}1
    {Î}{{\^I}}1 {Ô}{{\^O}}1 {Û}{{\^U}}1 {Ç}{{\c{C}}}1
    {²}{{\textsuperscript{2}}}1
    {°}{{$^\circ$}}1,
  morekeywords={String, ArrayList, HashSet, TreeSet, LinkedHashSet,
                HashMap, Hashtable, Iterator, Enumeration, Serializable,
                DataOutputStream, DataInputStream, FileOutputStream,
                FileInputStream, ObjectOutputStream, ObjectInputStream,
                BufferedReader, BufferedWriter, FileReader, FileWriter,
                PrintWriter, Scanner, Override, var, Date, synchronized,
                Runnable, Thread, wait, notifyAll, notify},
}
\lstset{style=java}

% ─── En-têtes et pieds de page ──────────────────────────────────────────────
\pagestyle{fancy}
\fancyhf{}
\fancyhead[L]{\small\color{bleuENSA}\textbf{ENSA Khouribga}%
  ~\textcolor{grisBorder}{$|$}~\textit{IID1~~·~~LABSER Yahya}}
\fancyhead[R]{\small\color{bleuAccent}\textsc{Module POO -- Java}}
\fancyfoot[L]{\small\color{grisTexte}\textit{Rapport de TPs -- 2024/2025}}
\fancyfoot[C]{\small\color{bleuENSA}\thepage}
\fancyfoot[R]{\small\color{grisTexte}\textit{Prof. N. GHERABI}}
\renewcommand{\headrulewidth}{0.6pt}
\renewcommand{\headrule}{%
  \hbox to\headwidth{\color{bleuENSA}\leaders\hrule height \headrulewidth\hfill}}
\renewcommand{\footrulewidth}{0.3pt}
\renewcommand{\footrule}{%
  \hbox to\headwidth{\color{grisBorder}\leaders\hrule height \footrulewidth\hfill}}

% ─── CORRECTION PRINCIPALE : Numérotation des TPs ──────────────────────────
\newcounter{tpnum}
\setcounter{tpnum}{0}

\renewcommand{\thesection}{\arabic{tpnum}.\arabic{section}}
\renewcommand{\thesubsection}{\arabic{tpnum}.\arabic{section}.\arabic{subsection}}

\newcommand{\tpchapter}[2]{%   #1 = numéro TP affiché, #2 = titre
  \clearpage
  \setcounter{tpnum}{#1}%
  \setcounter{section}{0}%
  \setcounter{subsection}{0}%
  \phantomsection
  \addcontentsline{toc}{chapter}{TP #1~~--~~#2}%
  \thispagestyle{fancy}%
  \noindent{\color{bleuAccent!40}\rule{\linewidth}{0.5pt}}\\[-6pt]
  \noindent{\color{bleuENSA}\rule{\linewidth}{3pt}}%
  \vspace{8pt}\par
  \noindent{%
    \parbox[t]{2.2cm}{%
      \color{bleuENSA}\LARGE\bfseries\scshape TP#1%
    }%
    \hspace{6pt}%
    \parbox[t]{\dimexpr\linewidth-2.8cm\relax}{%
      \color{bleuTitre}\LARGE\bfseries #2%
    }%
  }\par
  \vspace{6pt}%
  \noindent{\color{bleuENSA}\rule{\linewidth}{3pt}}\\[-4pt]
  \noindent{\color{bleuAccent!40}\rule{\linewidth}{0.5pt}}%
  \vspace{12pt}\par
}

% ─── Style des sections ─────────────────────────────────────────────────────
\titleformat{\section}
  {\color{bleuAccent}\large\bfseries}
  {\thesection}{0.8em}{}
  [\vspace{1pt}{\color{bleuAccent!35}\titlerule[0.6pt]}\vspace{4pt}]

\titleformat{\subsection}
  {\color{bleuENSA!80}\normalsize\bfseries}
  {\thesubsection}{0.7em}{}

\titlespacing{\section}{0pt}{14pt}{6pt}
\titlespacing{\subsection}{0pt}{10pt}{4pt}

% ─── Commandes screenshots ──────────────────────────────────────────────────
\newcommand{\screenshot}[2]{%
  \begin{center}
    \IfFileExists{#1}{%
      \fbox{\includegraphics[width=0.90\linewidth,keepaspectratio]{#1}}%
    }{%
      \begin{mdframed}[linecolor=grisBorder, linewidth=0.5pt,
        backgroundcolor=grisFond,
        innertopmargin=14pt, innerbottommargin=14pt,
        innerleftmargin=10pt, innerrightmargin=10pt]
        \centering\color{grisTexte}\small
        \faImage~~\textit{[Capture d'écran : \texttt{#1}]}
      \end{mdframed}%
    }%
    \par\vspace{3pt}
    {\footnotesize\color{grisTexte}\textit{#2}}%
  \end{center}%
  \vspace{4pt}%
}

% ─── Boîte Énoncé ───────────────────────────────────────────────────────────
\newmdenv[
  linecolor=sujBorder,
  linewidth=1pt,
  topline=true, rightline=true, bottomline=true, leftline=true,
  leftmargin=0pt,
  innerleftmargin=12pt, innertopmargin=8pt, innerbottommargin=8pt,
  innerrightmargin=12pt,
  backgroundcolor=sujBg,
  skipabove=10pt, skipbelow=10pt,
  roundcorner=3pt
]{enonce_env}

\newcommand{\enonce}[1]{%
  \begin{enonce_env}
    \textbf{\color{sujBorder}\faClipboardList~~Énoncé :}~#1
  \end{enonce_env}
}

% ─── Boîtes personnalisées ──────────────────────────────────────────────────
\newmdenv[
  linecolor=bleuENSA, linewidth=3pt,
  topline=false, rightline=false, bottomline=false,
  leftmargin=0pt, innerleftmargin=14pt,
  innertopmargin=10pt, innerbottommargin=10pt,
  backgroundcolor=bleuClair,
  skipabove=12pt, skipbelow=12pt, roundcorner=2pt
]{objectifBox}

\newmdenv[
  linecolor=violetInfo, linewidth=2.5pt,
  topline=false, rightline=false, bottomline=false,
  leftmargin=0pt, innerleftmargin=14pt,
  innertopmargin=9pt, innerbottommargin=9pt,
  backgroundcolor=violetClairBg,
  skipabove=10pt, skipbelow=10pt, roundcorner=2pt
]{observationBox}

\newmdenv[
  linecolor=orangeNote, linewidth=2.5pt,
  topline=false, rightline=false, bottomline=false,
  leftmargin=0pt, innerleftmargin=14pt,
  innertopmargin=9pt, innerbottommargin=9pt,
  backgroundcolor=orangeClairBg,
  skipabove=10pt, skipbelow=10pt, roundcorner=2pt
]{noteBox}

\newmdenv[
  linecolor=vertSucces, linewidth=2.5pt,
  topline=false, rightline=false, bottomline=false,
  leftmargin=0pt, innerleftmargin=14pt,
  innertopmargin=9pt, innerbottommargin=9pt,
  backgroundcolor=vertClairBg,
  skipabove=12pt, skipbelow=12pt, roundcorner=2pt
]{conclusionTPBox}

\newmdenv[
  linecolor=rougeErr, linewidth=2.5pt,
  topline=false, rightline=false, bottomline=false,
  leftmargin=0pt, innerleftmargin=14pt,
  innertopmargin=9pt, innerbottommargin=9pt,
  backgroundcolor=rougeClairBg,
  skipabove=10pt, skipbelow=10pt, roundcorner=2pt
]{erreurBox}

\newmdenv[
  linecolor=bleuMoyen, linewidth=2.5pt,
  topline=false, rightline=false, bottomline=false,
  leftmargin=0pt, innerleftmargin=14pt,
  innertopmargin=9pt, innerbottommargin=9pt,
  backgroundcolor=grisClairBg,
  skipabove=10pt, skipbelow=10pt, roundcorner=2pt
]{problemeBox}

% ─── Commandes raccourcis ────────────────────────────────────────────────────
\newcommand{\objectif}[1]{%
  \begin{objectifBox}
    \textbf{\color{bleuENSA}\faBullseye~~Objectif :}~~#1
  \end{objectifBox}
}
\newcommand{\observation}[1]{%
  \begin{observationBox}
    \textbf{\color{violetInfo}\faEye~~Observation :}~~#1
  \end{observationBox}
}
\newcommand{\note}[1]{%
  \begin{noteBox}
    \textbf{\color{orangeNote}\faInfoCircle~~Note :}~~#1
  \end{noteBox}
}
\newcommand{\conclusionTP}[1]{%
  \begin{conclusionTPBox}
    \textbf{\color{vertSucces}\faCheckCircle~~Conclusion :}~~#1
  \end{conclusionTPBox}
}
\newcommand{\erreurFrequente}[1]{%
  \begin{erreurBox}
    \textbf{\color{rougeErr}\faExclamationTriangle~~Erreur fréquente :}~~#1
  \end{erreurBox}
}
\newcommand{\probleme}[1]{%
  \begin{problemeBox}
    \textbf{\color{bleuMoyen}\faWrench~~Problème rencontré et solution :}~~#1
  \end{problemeBox}
}

% ─── Hyperlinks ─────────────────────────────────────────────────────────────
\hypersetup{
  colorlinks=true,
  linkcolor=bleuENSA,
  urlcolor=bleuAccent,
  bookmarksnumbered=true,
  bookmarksopen=true,
  pdftitle={Rapport TPs POO Java -- ENSA Khouribga -- LABSER Yahya},
  pdfauthor={LABSER Yahya},
  pdfsubject={Programmation Orientée Objet en Java},
}

% ════════════════════════════════════════════════════════════════════════════
\begin{document}

% ─── Page de garde ──────────────────────────────────────────────────────────
\begin{titlepage}
  \centering

  \begin{tikzpicture}[remember picture, overlay]
    \fill[bleuENSA] (current page.north west)
      rectangle ([yshift=-4.2cm]current page.north east);
    \fill[orDore] ([yshift=-4.2cm]current page.north west)
      rectangle ([yshift=-4.55cm]current page.north east);

    \node[anchor=north, text=white] at ([yshift=-0.75cm]current page.north)
      {\Large\bfseries\scshape École Nationale des Sciences Appliquées};
    \node[anchor=north, text=white!80] at ([yshift=-1.6cm]current page.north)
      {\large Khouribga~~·~~Université Sultan Moulay Slimane};
    \node[anchor=north, text=white!60] at ([yshift=-2.55cm]current page.north)
      {\normalsize Filière : Ingénierie Informatique et Données~|~IID1};
    \draw[white!40, line width=0.5pt]
      ([yshift=-3.4cm, xshift=3cm]current page.north west) --
      ([yshift=-3.4cm, xshift=-3cm]current page.north east);
  \end{tikzpicture}

  \vspace{5.0cm}

  {\color{bleuENSA!40}\rule{0.12\linewidth}{0.4pt}}~%
  {\color{orDore}\rule{0.60\linewidth}{1.8pt}}~%
  {\color{bleuENSA!40}\rule{0.12\linewidth}{0.4pt}}

  \vspace{0.7cm}

  {\fontsize{26}{32}\selectfont\bfseries\color{bleuENSA}
    Rapport de Travaux Pratiques}\\[0.4cm]
  {\LARGE\bfseries\color{grisTexte}
    Module : Programmation Orientée Objet}\\[0.3cm]
  {\Large\color{bleuMoyen}\itshape Langage Java}

  \vspace{0.6cm}

  {\color{bleuENSA!40}\rule{0.12\linewidth}{0.4pt}}~%
  {\color{orDore}\rule{0.60\linewidth}{1.8pt}}~%
  {\color{bleuENSA!40}\rule{0.12\linewidth}{0.4pt}}

  \vspace{1.2cm}

  \begin{mdframed}[linecolor=bleuENSA, linewidth=1.2pt,
    backgroundcolor=bleuClair,
    innertopmargin=14pt, innerbottommargin=14pt,
    innerleftmargin=26pt, innerrightmargin=26pt,
    roundcorner=5pt]
    \begin{tabular}{@{}>{\bfseries\color{bleuENSA}}l%
                       @{\hspace{1.2em}:\hspace{1.2em}}l@{}}
      Établissement & ENSA Khouribga \\[7pt]
      Filière       & IID1 \\[7pt]
      Étudiant      & \textbf{\Large LABSER Yahya} \\[7pt]
      Encadrant     & Prof.\ Noreddine GHERABI \\[7pt]
      Année         & 2024 -- 2025 \\
    \end{tabular}
  \end{mdframed}

  \vspace{0.8cm}

  \begin{mdframed}[linecolor=grisBorder, linewidth=0.8pt,
    backgroundcolor=white,
    innertopmargin=12pt, innerbottommargin=12pt,
    innerleftmargin=10pt, innerrightmargin=10pt,
    roundcorner=5pt]
    \centering
    {\large\bfseries\color{bleuENSA}Sommaire des Travaux Pratiques}\\[8pt]
    \renewcommand{\arraystretch}{1.25}
    \begin{tabularx}{0.97\linewidth}{@{}c X p{3.8cm}@{}}
      \toprule
      \rowcolor{bleuENSA!12}
      \textbf{TP} & \textbf{Thème} & \textbf{Concepts clés} \\
      \midrule
      1  & Concepts de base de la programmation Java    & Boucles, Switch, Classes \\
      2  & Tableaux, chaînes de caractères et méthodes  & Tableaux, \texttt{toString}, \texttt{equals} \\
      3  & Héritage simple -- Gestion immobilière       & \texttt{extends}, \texttt{super} \\
      4  & Classes abstraites et polymorphisme           & \texttt{abstract}, \texttt{@Override} \\
      5  & Collections Java -- Gestion des employés     & Set, Map, List, Iterator \\
      6  & Exceptions personnalisées -- Banque           & \texttt{try/catch}, \texttt{throw} \\
      7  & Threads et synchronisation -- DAB            & \texttt{synchronized}, \texttt{wait} \\
      9  & Gestion des fichiers et flux                 & IO, Sérialisation \\
      10 & Interfaces graphiques et base de données     & Swing, JDBC \\
      11 & Opérations CRUD avancées -- Swing \& JDBC   & INSERT, SELECT, UPDATE, JDBC \\
      \bottomrule
    \end{tabularx}
    \vspace{4pt}
    {\footnotesize\color{grisTexte}\textit{Note : Le TP~8 n'a pas été réalisé dans le cadre de ce module.}}
  \end{mdframed}

  \vfill
  {\small\color{gray}Document généré le \today}
\end{titlepage}

% ─── Table des matières ─────────────────────────────────────────────────────
\setcounter{tocdepth}{2}
\tableofcontents
\thispagestyle{empty}
\newpage

% ════════════════════════════════════════════════════════════════════════════
%  TP 1 : Concepts de base
% ════════════════════════════════════════════════════════════════════════════
\tpchapter{1}{Concepts de Base de la Programmation Java}

\objectif{Maîtriser les instructions fondamentales de Java : boucles
(\texttt{while}, \texttt{for}), instructions de contrôle (\texttt{continue},
\texttt{switch}), méthodes et définition de la première classe orientée objet.}

% ── Exercice 1 ──────────────────────────────────────────────────────────────
\section{Exercice 1 -- Somme des parties entières}

\enonce{Écrire un programme qui demande à l'utilisateur une série de nombres
réels (arrêt sur 0) et affiche la somme des parties entières.}

\begin{lstlisting}[caption={\texttt{TesterNbr.java} -- Somme des parties entières}]
package TP1.ex1;
import java.util.Scanner;

public class TesterNbr {
    public static void main(String[] args) {
        int s = 0;
        Scanner scan = new Scanner(System.in);
        String line = scan.nextLine();
        while (!line.equals("0")) {
            s = s + Integer.parseInt(line);
            line = scan.nextLine();
        }
        System.out.println(s);
    }
}
\end{lstlisting}

\screenshot{Screenshot 2026-05-03 155327.png}{Code Java -- Somme des parties entières et résultat d'exécution}

\observation{La méthode \texttt{Integer.parseInt()} tronque la partie décimale
sans arrondir. Un appel avec \texttt{"3.7"} lève une
\texttt{NumberFormatException} ; il faut d'abord passer par
\texttt{Double.parseDouble()} puis caster en \texttt{int} pour extraire
correctement la partie entière.}

\note{La condition d'arrêt \texttt{while(!line.equals("0"))} suppose une saisie
exactement égale à \texttt{"0"}. Une amélioration serait d'utiliser
\texttt{Double.parseDouble(line) == 0} pour accepter \texttt{"0.0"},
\texttt{"0.00"}, etc.}

\probleme{Lors de la première version, \texttt{Integer.parseInt()} était utilisé
directement sur la saisie, levant une \texttt{NumberFormatException} pour tout
nombre décimal. La solution a été de passer d'abord par \texttt{Double.parseDouble()},
puis de caster en \texttt{int} pour obtenir la partie entière sans erreur.}

% ── Exercice 2 ──────────────────────────────────────────────────────────────
\section{Exercice 2 -- Carrés avec \texttt{continue}}

\enonce{Afficher le carré des entiers de $-3$ à $+3$ en utilisant
l'instruction \texttt{continue}.}

\begin{lstlisting}[caption={\texttt{Carre.java} -- Affichage des carrés avec \texttt{continue}}]
package TP1.ex2;

public class Carre {
    public static void main(String[] args) {
        int n = -3;
        while (n <= 3) {
            if (n == 0) {
                n++;
                continue;
            }
            System.out.println(n * n);
            n++;
        }
    }
}
\end{lstlisting}

\screenshot{Screenshot 2026-05-03 160201.png}{Code Java -- Carrés avec \texttt{continue} et résultat d'exécution}

\observation{L'instruction \texttt{continue} interrompt l'itération courante et
passe directement à la suivante. Elle exclut ici $n = 0$ du calcul du carré,
illustrant son rôle dans le contrôle de flux sans bloc \texttt{if/else} imbriqué.}

\note{Le carré de $0$ vaut $0$ -- son exclusion n'a aucun impact mathématique
ici. L'exercice vise uniquement à pratiquer syntaxiquement le mot-clé
\texttt{continue}.}

% ── Exercice 3 ──────────────────────────────────────────────────────────────
\section{Exercice 3 -- Racines carrées des entiers pairs}

\enonce{Afficher la racine carrée des entiers pairs de 0 à 10 inclus.}

\begin{lstlisting}[caption={\texttt{Racine.java} -- Racines carrées des entiers pairs}]
package TP1.ex3;

public class Racine {
    public static void main(String[] args) {
        for (int i = 0; i <= 10; i += 2) {
            System.out.println("la racine carre de " + i + " est: " + Math.sqrt(i));
        }
    }
}
\end{lstlisting}

\screenshot{Screenshot 2026-05-03 160302.png}{Code Java -- Racines carrées des pairs et résultat d'exécution}

\observation{L'incrémentation \texttt{i += 2} garantit de ne parcourir que les
entiers pairs, plus efficace qu'un test \texttt{if (i \% 2 == 0)} à chaque
itération. \texttt{Math.sqrt()} retourne un \texttt{double}, adapté aux
résultats irrationnels comme $\sqrt{2} \approx 1{,}414$.}

% ── Exercice 4 ──────────────────────────────────────────────────────────────
\section{Exercice 4 -- Factorielle}

\enonce{Calculer la factorielle d'un entier saisi par l'utilisateur et
afficher le résultat.}

\begin{lstlisting}[caption={\texttt{Fact.java} -- Calcul de la factorielle}]
package TP1.ex4;
import java.util.Scanner;

public class Fact {
    public static void main(String[] args) {
        Scanner scan = new Scanner(System.in);
        System.out.print("Entrez un entier : ");
        int nbr = Integer.parseInt(scan.nextLine());
        if (nbr < 0)
            System.exit(-1);
        int result = 1;
        for (int i = 1; i <= nbr; i++)
            result *= i;
        System.out.println(result);
    }
}
\end{lstlisting}

\screenshot{Screenshot 2026-05-03 160408.png}{Code Java -- Calcul de la factorielle et résultat d'exécution}

\observation{Le type \texttt{int} (32 bits) déborde à partir de $n \geq 13$ :
$13! = 6\,227\,020\,800$ dépasse \texttt{Integer.MAX\_VALUE}. Pour de grandes
valeurs, il faut recourir à \texttt{long} ou \texttt{BigInteger}.}

\erreurFrequente{Initialiser \texttt{result = 0} au lieu de \texttt{result = 1}
retourne systématiquement $0$. L'élément neutre de la multiplication est $1$,
pas $0$.}

\probleme{L'accumulateur initialisé à \texttt{0} donnait toujours $0$ en sortie.
En comprenant que la multiplication nécessite l'élément neutre $1$ pour démarrer,
l'initialisation a été corrigée et le programme a fonctionné correctement.}

% ── Exercice 5 ──────────────────────────────────────────────────────────────
\section{Exercice 5 -- Calculatrice avec \texttt{switch}}

\enonce{En utilisant l'instruction de choix multiples \texttt{switch}, écrire un
programme qui initialise trois valeurs et calcule le résultat selon
l'opérateur \texttt{(+, -, *)} saisi par l'utilisateur.}

\begin{lstlisting}[caption={\texttt{Calcule.java} -- Calculatrice utilisant \texttt{switch} avec flèches}]
package TP1.ex5;
import java.util.Scanner;

public class Calcule {
    public static void main(String[] args) {
        Scanner scan = new Scanner(System.in);
        System.out.println("Entre a: ");
        int a = Integer.parseInt(scan.nextLine());
        System.out.println("Entre b: ");
        int b = Integer.parseInt(scan.nextLine());
        System.out.println("Entre c: ");
        int c = Integer.parseInt(scan.nextLine());

        System.out.println("Entre l'operateur: ");
        String operator = scan.nextLine();

        int result;
        switch (operator) {
            case "+" -> result = a + b + c;
            case "-" -> result = a - b - c;
            case "*" -> result = a * b * c;
            default -> {
                System.out.println("opreation n'est pas connais!");
                return;
            }
        };
        System.out.println("la resultat est : " + result);
    }
}
\end{lstlisting}

\screenshot{Screenshot 2026-05-03 160521.png}{Code Java -- Calculatrice avec \texttt{switch} et résultat d'exécution}

\observation{Java 14+ introduit la syntaxe \texttt{switch} avec les flèches
(\texttt{->}), supprimant la nécessité du \texttt{break} et améliorant la
lisibilité. Le code présenté utilise déjà cette syntaxe moderne.}

\note{La division par zéro n'est pas gérée dans cet exercice. En situation
réelle, un test \texttt{if (b == 0 \&\& operator.equals("/"))} devrait être
ajouté avant toute division pour éviter une \texttt{ArithmeticException}.}

\probleme{Lors des premiers tests, une division par zéro levait une
\texttt{ArithmeticException} sans message explicite. Après ajout d'une
vérification préalable de l'opérande, un message métier clair est affiché à
l'utilisateur.}

% ── Exercice 6 ──────────────────────────────────────────────────────────────
\section{Exercice 6 -- Classe \texttt{Produit}}

\enonce{Implémenter la classe \texttt{Produit} avec les attributs
\texttt{Code}, \texttt{Intitulé}, \texttt{Prix} et \texttt{Marque}, deux
constructeurs, des accesseurs, une méthode \texttt{afficher()}, une méthode
\texttt{augmenter(M)} (M $\leq$ 5\,\% du prix) et une méthode
\texttt{diminuer(N\%)} (uniquement si prix $> 1\,000$ DH). Créer deux
instances P1 et P2 dans un programme principal.}

\screenshot{Screenshot 2026-05-03 160731.png}{Classe \texttt{Produit} -- attributs et constructeurs}
\screenshot{Screenshot 2026-05-03 160745.png}{Classe \texttt{Produit} -- méthodes \texttt{augmenter} et \texttt{diminuer}}
\screenshot{Screenshot 2026-05-03 160756.png}{Programme principal -- création de P1, P2 et résultat d'exécution}

\begin{lstlisting}[caption={\texttt{Produit.java} -- attributs, constructeurs et méthodes métier}]
package TP1.ex6;

public class Produit {
    int code;
    String intitule;
    float prix;
    String marque;

    public Produit() {
        code = 0;
        intitule = "Default";
        prix = 0;
        marque = "Default";
    }

    public Produit(int code, float prix) {
        this.code = code;
        this.intitule = "Default";
        this.prix = prix;
        this.marque = "Default";
    }

    public float getPrix()    { return prix; }
    public String getMarque() { return marque; }

    public void afficher() {
        System.out.println("Code: " + code);
        System.out.println("Intitule: " + intitule);
        System.out.println("Prix: " + getPrix());
        System.out.println("Marque: " + getMarque());
    }

    public void augmenter(float m) {
        if (m <= 0.05 * prix)
            prix = prix + m;
        else
            System.out.println("Montant dépasse 5% du prix !");
    }

    public void diminuer(float n) {
        if (prix > 1000)
            prix = prix - (n / 100 * prix);
        else
            System.out.println("Prix doit être supérieur à 1000dh !");
    }
}
\end{lstlisting}

\begin{lstlisting}[caption={\texttt{TestProduit.java} -- création de P1, P2 et appels métier}]
package TP1.ex6;

public class TestProduit {
    public static void main(String[] args) {
        Produit p1 = new Produit(1, 30000);
        Produit p2 = new Produit(2, 3500);

        p1.augmenter(300);
        p2.diminuer(3);

        p1.afficher();
        p2.afficher();
    }
}
\end{lstlisting}

\observation{La classe \texttt{Produit} illustre le principe
d'\textbf{encapsulation} : les méthodes \texttt{augmenter} et \texttt{diminuer}
contrôlent les modifications selon des règles métier précises -- augmentation
limitée à 5\,\%, prix minimum de 1\,000 DH pour la réduction -- empêchant
toute modification incohérente depuis l'extérieur.}

\conclusionTP{Ce premier TP a consolidé les bases essentielles de Java : syntaxe
des boucles (\texttt{while}, \texttt{for}), instructions de contrôle
(\texttt{continue}, \texttt{switch}), calculs numériques, et première définition
d'une classe avec constructeurs, accesseurs et méthodes métier. Ces fondamentaux
sont indispensables pour aborder la programmation orientée objet.}

% ════════════════════════════════════════════════════════════════════════════
%  TP 2 : Tableaux, chaînes, méthodes
% ════════════════════════════════════════════════════════════════════════════
\tpchapter{2}{Tableaux, Chaînes de Caractères et Méthodes}

\objectif{Maîtriser la gestion des tableaux à taille fixe, la manipulation
des chaînes de caractères (\texttt{toUpperCase}, \texttt{replace}), et la
définition de méthodes statiques utilitaires.}

\section{Exercice 1 -- Traitement de texte avec \texttt{switch}}

\enonce{Demander à l'utilisateur un texte T et un mot A. Utiliser
\texttt{switch} pour : supprimer les occurrences d'\texttt{"UNIVERSITE"},
remplacer \texttt{"FACULTE"} par \texttt{"UNIVERSITE"}, insérer
\texttt{"ECOLE"} à la fin, ou afficher \textit{"Opération non reconnue"}.}

\begin{lstlisting}[caption={\texttt{TexteOperation.java} -- traitement de texte avec \texttt{switch}}]
package TP2.ex1;
import java.util.Scanner;

public class TexteOperation {
    public static void main(String[] args) {
        Scanner scan = new Scanner(System.in);
        System.out.println("Entrer un texte : ");
        String t = scan.nextLine();
        System.out.println("Entrer un mot : ");
        String a = scan.nextLine().toUpperCase();

        switch (a) {
            case "UNIVERSITE" -> t = t.replace("UNIVERSITE", "");
            case "FACULTE"    -> t = t.replace("FACULTE", "UNIVERSITE");
            case "ECOLE"      -> t = t + " ECOLE";
            default           -> {
                System.out.println("Opération non reconnue.");
                return;
            }
        }
        System.out.println(t);
    }
}
\end{lstlisting}

\screenshot{Screenshot 2026-05-03 161020.png}{Code -- Traitement de texte avec \texttt{switch} et résultat d'exécution}

\observation{L'appel \texttt{toUpperCase()} sur la saisie \texttt{a} normalise
la comparaison dans le \texttt{switch}, rendant le programme insensible à la
casse. \texttt{String.replace()} retourne une \emph{nouvelle} chaîne (les
\texttt{String} Java sont immuables) -- le résultat doit être réassigné à
\texttt{t}.}

\probleme{La première version était sensible à la casse : \texttt{"majuscule"}
ne correspondait pas à \texttt{"MAJUSCULE"}. En appliquant \texttt{toUpperCase()}
avant la comparaison dans le \texttt{switch}, le problème a été résolu de façon
élégante.}

\section{Exercice 2 -- Tableau d'entiers}

\enonce{Écrire une méthode statique \texttt{afficherTableau} prenant un
tableau d'entiers en paramètre et affichant ses éléments. Dans
\texttt{main}, demander la taille, saisir les valeurs, puis appeler la méthode.}

\begin{lstlisting}[caption={\texttt{Tableau.java} -- méthode statique \texttt{afficherTableau} et saisie d'un tableau}]
package TP2.ex2;
import java.util.Scanner;

public class Tableau {
    static void afficherTableau(int[] tab) {
        int N = tab.length;
        for (int i = 0; i < N; i++) {
            System.out.println("La valeur N" + i + " est: " + tab[i]);
        }
    }

    public static void main(String[] args) {
        Scanner scan = new Scanner(System.in);
        System.out.println("Entre la taille du tableau: ");
        int N = Integer.parseInt(scan.nextLine());

        int[] tab = new int[N];
        for (int i = 0; i < N; i++) {
            System.out.println("Entree la valeur N" + i);
            tab[i] = Integer.parseInt(scan.nextLine());
        }
        afficherTableau(tab);
    }
}
\end{lstlisting}

\screenshot{Screenshot 2026-05-03 161122.png}{Code -- Saisie et affichage d'un tableau via méthode statique}

\observation{Déclarer \texttt{afficherTableau} comme \texttt{static} permet
de l'appeler sans instancier la classe, pratique recommandée pour les méthodes
utilitaires ne dépendant d'aucun état d'objet.}

\note{Les tableaux Java ont une taille fixe définie à la création. Pour des
collections de taille variable, \texttt{ArrayList<Integer>} est préférable.}

\section{Exercice 3 -- Classe \texttt{Etudiant}}

\enonce{Implémenter la classe \texttt{Etudiant} (attributs privés
\texttt{id}, \texttt{nom}, \texttt{note} ; deux constructeurs ; getters/setters ;
\texttt{toString()} ; \texttt{equals()} basé sur \texttt{id + nom + note}).
Dans \texttt{main}, saisir $n$ étudiants, calculer la moyenne via
\texttt{calculerMoyenne()}, et comparer deux étudiants avec \texttt{equals()}.}

\screenshot{Screenshot 2026-05-03 161348.png}{Classe \texttt{Etudiant} -- définition, \texttt{toString} et \texttt{equals}}
\screenshot{Screenshot 2026-05-03 161439.png}{Programme principal -- saisie, moyenne et comparaison d'étudiants}

\begin{lstlisting}[caption={\texttt{Etudiant.java} -- classe avec \texttt{toString} et \texttt{equals}}]
package TP2.ex3;

public class Etudiant {
    private int id;
    private String nom;
    private float note;

    public Etudiant() {
        id = 0; nom = "Default"; note = 0;
    }
    public Etudiant(int id, String nom, float note) {
        this.id = id; this.nom = nom; this.note = note;
    }

    public int getId()      { return id; }
    public String getNom()  { return nom; }
    public float getNote()  { return note; }

    public void setId(int id)        { this.id = id; }
    public void setNom(String nom)   { this.nom = nom; }
    public void setNote(float note)  { this.note = note; }

    public String toString() {
        return "id: " + id + ", nom: " + nom + ", note: " + note;
    }

    public boolean equals(Etudiant other) {
        return this.id == other.id
            && this.nom.equals(other.nom)
            && this.note == other.note;
    }
}
\end{lstlisting}

\begin{lstlisting}[caption={\texttt{TestEtudiant.java} -- saisie, moyenne et comparaison}]
package TP2.ex3;
import java.util.Scanner;

public class TestEtudiant {
    static float calculerMoyenne(Etudiant[] tab) {
        float s = 0;
        for (int i = 0; i < tab.length; i++)
            s += tab[i].getNote();
        return s / tab.length;
    }

    public static void main(String[] args) {
        Scanner scan = new Scanner(System.in);
        System.out.println("Entree le nbr des etudiants: ");
        int N = Integer.parseInt(scan.nextLine());

        Etudiant[] tab = new Etudiant[N];
        for (int i = 0; i < N; i++) {
            tab[i] = new Etudiant();
            System.out.println("Entre l'id d'etudiant " + i);
            tab[i].setId(Integer.parseInt(scan.nextLine()));
            System.out.println("Entree le nom d'etudiant " + i);
            tab[i].setNom(scan.nextLine());
            System.out.println("Entree la note d'etudiant " + i);
            tab[i].setNote(Float.parseFloat(scan.nextLine()));
        }

        System.out.println("Les etudiants:");
        for (int i = 0; i < N; i++)
            System.out.println(tab[i].toString());

        System.out.println("La moyenne des notes: " + calculerMoyenne(tab));

        System.out.println("Entre deux indices i et j :");
        int i = Integer.parseInt(scan.nextLine());
        int j = Integer.parseInt(scan.nextLine());

        if (i >= 0 && j >= 0 && i < N && j < N) {
            if (tab[i].equals(tab[j]))
                System.out.println("Les deux etudiants sont identiques");
            else
                System.out.println("Les deux etudiants ne sont pas identiques");
        } else
            System.out.println("Indices invalides ! Doit être entre 0 et " + (N - 1));
    }
}
\end{lstlisting}

\observation{La redéfinition de \texttt{toString()} permet l'affichage direct
via \texttt{System.out.println(objet)}. En pratique, redéfinir aussi
\texttt{hashCode()} est recommandé pour respecter le contrat général de Java
lorsqu'on redéfinit \texttt{equals()}.}

\erreurFrequente{Utiliser \texttt{==} pour comparer deux \texttt{String} compare
les références mémoire, pas le contenu. Il faut toujours utiliser \texttt{.equals()}
pour les comparaisons de chaînes.}

\probleme{Dans \texttt{equals()}, les champs \texttt{String} étaient comparés
avec \texttt{==}, retournant \texttt{false} même pour des étudiants identiques.
Remplacer \texttt{==} par \texttt{.equals()} a immédiatement corrigé l'erreur.}

\conclusionTP{Ce TP a consolidé la manipulation des tableaux à taille fixe, les
méthodes statiques utilitaires, et l'enrichissement d'une classe via
\texttt{toString()} et \texttt{equals()}. La notion d'immuabilité des
\texttt{String} en Java a été mise en pratique concrètement.}

% ════════════════════════════════════════════════════════════════════════════
%  TP 3 : Héritage simple
% ════════════════════════════════════════════════════════════════════════════
\tpchapter{3}{Héritage Simple -- Gestion Immobilière}

\objectif{Définir une hiérarchie de classes par héritage simple, utiliser
\texttt{super()} pour appeler le constructeur parent, surcharger
\texttt{calculeValeurLocative()} dans chaque sous-classe selon les règles
métier spécifiées.}

\section{Classe de base \texttt{BienImmobilier}}

\enonce{Définir la classe \texttt{BienImmobilier} avec les attributs
\texttt{adresse}, \texttt{surface}, \texttt{prixAchat} et
\texttt{valeurLocativeEstimee}. Fournir un constructeur (3 paramètres),
une méthode \texttt{affiche()}, et une méthode
\texttt{calculeValeurLocative()} fixant la valeur à $0{,}5\,\%$ du prix
d'achat par m².}

\screenshot{Screenshot 2026-05-03 161643.png}{Classe \texttt{BienImmobilier} -- attributs, constructeur et \texttt{calculeValeurLocative}}

\begin{lstlisting}[caption={\texttt{BienImmobilier.java} -- classe de base de la hiérarchie}]
package TP3;

public class BienImmobilier {
    String adresse;
    float surface;
    float prixAcha;
    float valeurEst;

    public BienImmobilier() {
        adresse = "Deafault";
        surface = 0;
        prixAcha = 0;
        valeurEst = 0;
    }

    public BienImmobilier(String adress, float surface, float prixAcha) {
        this.adresse = adress;
        this.surface = surface;
        this.prixAcha = prixAcha;
    }

    public void affiche() {
        System.out.println("Adress: " + adresse);
        System.out.println("Surface: " + surface);
        System.out.println("Prix d'achat: " + prixAcha);
        System.out.println("La valeur locative estimée: " + valeurEst);
    }

    public void calculeValeurLocative() {
        valeurEst = (float) 0.005 * prixAcha;
    }
}
\end{lstlisting}

\observation{La classe \texttt{BienImmobilier} définit le comportement commun à
tous les biens. La méthode \texttt{calculeValeurLocative()} est conçue pour être
redéfinie (\textit{overridden}) par les sous-classes, posant les fondations du
polymorphisme.}

\section{Classe \texttt{Maison}}

\enonce{Classe héritant de \texttt{BienImmobilier} avec \texttt{nbPieces},
\texttt{surfaceJardin} et \texttt{piscine}. Redéfinir
\texttt{calculeValeurLocative()} : base $= 0{,}5\,\% \times$ prix/m²,
$+10\,\%$ si jardin $> 50$ m², $+20\,\%$ si piscine.}

\screenshot{Screenshot 2026-05-03 161714.png}{Classe \texttt{Maison} héritant de \texttt{BienImmobilier} et redéfinissant \texttt{calculeValeurLocative}}

\begin{lstlisting}[caption={\texttt{Maison.java} -- héritage et surcharge de \texttt{calculeValeurLocative}}]
package TP3;

public class Maison extends BienImmobilier {
    int nbrP;
    float surfaceJ;
    boolean existP;

    public Maison() {
        nbrP = 0; surfaceJ = 0; existP = false;
    }

    public Maison(int nbrP, float surfaceJ, boolean existP) {
        super();
        this.nbrP = nbrP;
        this.surfaceJ = surfaceJ;
        this.existP = existP;
    }

    public void affiche() {
        super.affiche();
        System.out.println("Nombre de pièces: " + nbrP);
        System.out.println("Surface du jardin: " + surfaceJ);
        System.out.println("Exist une piscine: " + existP);
    }

    public void calculeValeurLocative() {
        super.calculeValeurLocative();
        if (surfaceJ > 50) valeurEst *= (float) 1.1;
        if (existP)        valeurEst *= (float) 1.2;
    }
}
\end{lstlisting}

\observation{La classe \texttt{Maison} appelle
\texttt{super.calculeValeurLocative()} avant d'appliquer ses facteurs
multiplicatifs (jardin $\times 1{,}1$, piscine $\times 1{,}2$). Ce chaînage
évite la duplication de code et respecte le principe \textit{Open/Closed}.}

\section{Classe \texttt{LocalCommercial}}

\enonce{Classe héritant de \texttt{BienImmobilier} avec \texttt{typeCommerce} ,
\texttt{locationMeublee} et \texttt{loyerMensuel}. Redéfinir
\texttt{calculeValeurLocative()} : base $= 0{,}8\,\% \times$ prix/m²,
$+15\,\%$ si meublé, $+5\,\%$ si Boutique, $-10\,\%$ si Entrepôt.
Résultat $\geq 0$.}

\screenshot{Screenshot 2026-05-03 161802.png}{Classe \texttt{LocalCommercial} -- calcul spécifique de la valeur locative}

\begin{lstlisting}[caption={\texttt{LocalCommercial.java} -- héritage avec règles métier différentes}]
package TP3;

public class LocalCommercial extends BienImmobilier {
    String typeC;
    boolean locationM;
    float loyerM;

    public LocalCommercial() {
        typeC = "Default"; locationM = false; loyerM = 0;
    }

    public LocalCommercial(String typeC, boolean locationM, float loyerM) {
        super();
        this.typeC = typeC;
        this.locationM = locationM;
        this.loyerM = loyerM;
    }

    public void affiche() {
        super.affiche();
        System.out.println("Type de commerce: " + typeC);
        System.out.println("La location meublée: " + locationM);
        System.out.println("Loyer mensuel actuel: " + loyerM);
    }

    public void calculeValeurLocative() {
        valeurEst = prixAcha * (float) 0.008;
        if (locationM) valeurEst *= (float) 1.15;
        switch (typeC) {
            case "Boutique" -> valeurEst *= (float) 1.05;
            case "Entrepôt" -> valeurEst *= (float) 0.9;
        }
        if (valeurEst < 0) valeurEst = 0;
    }
}
\end{lstlisting}

\observation{\texttt{LocalCommercial} recalcule entièrement sa valeur avec un
taux différent (0,8\,\% vs 0,5\,\%). La garde \texttt{if (valeurEst < 0)}
protège contre un résultat négatif causé par des combinaisons de coefficients
défavorables, conformément à la contrainte du sujet.}

\section{Classe de test \texttt{TestImmobilier}}

\enonce{Demander à l'utilisateur les informations pour une maison et un local
commercial, créer les objets, calculer leur valeur locative et afficher leur
état complet via \texttt{affiche()}.}

\screenshot{Screenshot 2026-05-03 161821.png}{Classe \texttt{TestImmobilier} -- saisie utilisateur, calcul et affichage}

\begin{lstlisting}[caption={\texttt{TestImmobilier.java} -- programme principal testant la hiérarchie}]
package TP3;
import java.util.Scanner;

public class TestImmobilier {
    public static void main(String[] args) {
        Scanner scan = new Scanner(System.in);

        System.out.println("Entre les info de la maison:");
        Maison m = new Maison();
        m.adresse  = scan.nextLine();
        m.surface  = Float.parseFloat(scan.nextLine());
        m.prixAcha = Float.parseFloat(scan.nextLine());
        m.nbrP     = Integer.parseInt(scan.nextLine());
        m.existP   = Boolean.parseBoolean(scan.nextLine());
        m.surfaceJ = Float.parseFloat(scan.nextLine());

        System.out.println("Entre les info de Local Commercial:");
        LocalCommercial l = new LocalCommercial();
        l.adresse   = scan.nextLine();
        l.surface   = Float.parseFloat(scan.nextLine());
        l.prixAcha  = Float.parseFloat(scan.nextLine());
        l.typeC     = scan.nextLine();
        l.locationM = Boolean.parseBoolean(scan.nextLine());
        l.loyerM    = Float.parseFloat(scan.nextLine());

        m.calculeValeurLocative();
        l.calculeValeurLocative();

        m.affiche();
        l.affiche();
    }
}
\end{lstlisting}

\note{Dans le programme de test, les attributs sont accédés directement (sans
setter) car leur modificateur d'accès est package-private par défaut. La bonne
pratique serait de les déclarer \texttt{private} et d'utiliser des constructeurs
paramétrés complets avec setters.}

\probleme{En omettant l'appel \texttt{super(...)} dans les constructeurs des
sous-classes, la compilation échouait car la classe parente ne possédait pas
de constructeur sans argument. En ajoutant systématiquement
\texttt{super(...)} en première instruction de chaque constructeur fils, le
problème a été résolu.}

\conclusionTP{Ce TP a mis en pratique le mécanisme fondamental de l'héritage :
la relation «~est-un~», l'appel \texttt{super()}, et la redéfinition de
méthodes. La hiérarchie
\texttt{BienImmobilier} $\to$ \texttt{Maison} / \texttt{LocalCommercial}
démontre comment factoriser le comportement commun tout en spécialisant le
comportement propre à chaque type de bien.}

% ════════════════════════════════════════════════════════════════════════════
%  TP 4 : Classes abstraites et polymorphisme
% ════════════════════════════════════════════════════════════════════════════
\tpchapter{4}{Classes Abstraites et Polymorphisme -- Gestion de Flotte}

\objectif{Définir une classe abstraite \texttt{Vehicule} avec méthodes
abstraites et concrètes, créer deux classes concrètes
(\texttt{VehiculeThermique}, \texttt{VehiculeElectrique}), et exploiter
le polymorphisme par substitution dans une liste unique de véhicules.}

\section{Classe abstraite \texttt{Vehicule}}

\enonce{Créer la classe abstraite \texttt{Vehicule} avec les attributs
privés \texttt{immatriculation}, \texttt{capacite} et
\texttt{consommationBase}, un constructeur paramétré, les méthodes abstraites
\texttt{calculerAutonomie()} et \texttt{calculerEfficacite()}, et les
getters/setters.}

\screenshot{Screenshot 2026-05-03 162150.png}{Classe abstraite \texttt{Vehicule} -- attributs, constructeurs et méthodes abstraites}

\begin{lstlisting}[caption={\texttt{Vehicule.java} -- classe abstraite définissant le contrat commun}]
package TP4;

public abstract class Vehicule {
    private String immatriculation;
    private int capacite;
    private double consommationBase;

    public Vehicule() {
        immatriculation = "default";
        capacite = 0;
        consommationBase = 0;
    }

    public Vehicule(String immatriculation, int capacite, double consommationBase) {
        this.immatriculation = immatriculation;
        this.capacite = capacite;
        this.consommationBase = consommationBase;
    }

    public String getImmatriculation()    { return immatriculation; }
    public int    getCapacite()           { return capacite; }
    public double getConsommationBase()   { return consommationBase; }

    public void setImmatriculation(String immatriculation) { this.immatriculation = immatriculation; }
    public void setCapacite(int capacite)                  { this.capacite = capacite; }
    public void setConsommationBase(double v)              { this.consommationBase = v; }

    public abstract double calculerAutonomie();
    public abstract double calculerEfficacite();
}
\end{lstlisting}

\observation{Déclarer \texttt{Vehicule} comme \texttt{abstract} empêche son
instanciation directe et impose à toute sous-classe concrète d'implémenter
\texttt{calculerAutonomie()} et \texttt{calculerEfficacite()}. Les champs
\texttt{private} avec accesseurs respectent l'encapsulation.}

\section{Classe \texttt{VehiculeThermique}}

\enonce{Créer \texttt{VehiculeThermique} héritant de \texttt{Vehicule},
avec \texttt{niveauCarburant}, \texttt{typeCarburant} et \texttt{cylindree}.
Formules : Autonomie $= \frac{\text{niveauCarburant} \times 100}{\text{consommationBase}}$,
Efficacité $= \frac{\text{capacité} \times 100}{\text{consommationBase} \times \text{cylindree}/1000}$,
Pollution $= \text{consommationBase} \times \text{cylindree} \times 0{,}02$.}

\screenshot{Screenshot 2026-05-03 162227.png}{Classe \texttt{VehiculeThermique} -- implémentation des méthodes et \texttt{calculerPollution}}

\begin{lstlisting}[caption={\texttt{VehiculeThermique.java} -- sous-classe concrète avec méthode propre}]
package TP4;

public class VehiculeThermique extends Vehicule {
    private double niveauCarburant;
    private String typeCarburant;
    private int cylindree;

    public VehiculeThermique() {
        super();
        niveauCarburant = 0;
        typeCarburant = "Default";
        cylindree = 0;
    }

    public double getNiveauCarburant()        { return niveauCarburant; }
    public String getTypeCarburant()          { return typeCarburant; }
    public int    getCylindree()              { return cylindree; }
    public void setNiveauCarburant(double v)  { this.niveauCarburant = v; }
    public void setTypeCarburant(String v)    { this.typeCarburant = v; }
    public void setCylindree(int v)           { this.cylindree = v; }

    public double calculerAutonomie() {
        return niveauCarburant * 100 / getConsommationBase();
    }
    public double calculerEfficacite() {
        return (getCapacite() * 100) / (getConsommationBase() * cylindree / 1000);
    }
    public double calculerPollution() {
        return getConsommationBase() * cylindree * 0.02;
    }

    @Override
    public String toString() {
        return ("immatriculation: " + this.getImmatriculation() +
                "\nautonomie,: "    + this.calculerAutonomie() +
                "\nefficacité: "    + this.calculerEfficacite() +
                "\npollution: "     + this.calculerPollution());
    }
}
\end{lstlisting}

\observation{La méthode \texttt{calculerPollution()} est propre à
\texttt{VehiculeThermique} et n'appartient pas au contrat abstrait. Cela montre
qu'une sous-classe peut ajouter ses propres comportements spécialisés
au-delà du contrat imposé par la classe parente.}

\section{Classe \texttt{VehiculeElectrique}}

\enonce{Créer \texttt{VehiculeElectrique} héritant de \texttt{Vehicule},
avec \texttt{niveauBatterie}, \texttt{tempsRecharge} et
\texttt{capaciteBatterie}. Formules : Autonomie $= \frac{\text{niveauBatterie} \times 100}{\text{consommationBase} \times 0{,}8}$,
Efficacité $= \frac{\text{capacité} \times \text{autonomie}}{\text{capacitéBatterie} \times \text{tempsRecharge}/60}$,
TempsRechargeRéel $= \text{tempsRecharge} \times (1 - \frac{\text{niveauBatterie}}{\text{capacitéBatterie}})$.}

\screenshot{Screenshot 2026-05-03 162328.png}{Classe \texttt{VehiculeElectrique} -- implémentation des méthodes et \texttt{calculerTempsRechargeReel}}

\begin{lstlisting}[caption={\texttt{VehiculeElectrique.java} -- sous-classe avec méthode \texttt{calculerTempsRechargeReel}}]
package TP4;

public class VehiculeElectrique extends Vehicule {
    private double niveauBatterie;
    private int tempsRecharge;
    private int capaciteBatterie;

    public VehiculeElectrique() {
        super();
        niveauBatterie = 0;
        tempsRecharge = 0;
        capaciteBatterie = 0;
    }

    public VehiculeElectrique(String immatriculation, int capacite, double consommationBase,
                              double niveauBatterie, int tempsRecharge, int capaciteBatterie) {
        super(immatriculation, capacite, consommationBase);
        this.niveauBatterie = niveauBatterie;
        this.tempsRecharge = tempsRecharge;
        this.capaciteBatterie = capaciteBatterie;
    }

    public double calculerAutonomie() {
        return niveauBatterie * 100 / (getConsommationBase() * 0.8);
    }
    public double calculerEfficacite() {
        return (getCapacite() * calculerAutonomie())
             / (capaciteBatterie * tempsRecharge / 60);
    }
    public double calculerTempsRechargeReel() {
        return tempsRecharge * (1 - niveauBatterie / capaciteBatterie);
    }

    @Override
    public String toString() {
        return ("immatriculation: " + this.getImmatriculation() +
                "\nautonomie,: "    + this.calculerAutonomie() +
                "\nefficacité: "    + this.calculerEfficacite() +
                "\ntemps de recharge réel: " + this.calculerTempsRechargeReel());
    }
}
\end{lstlisting}

\observation{\texttt{VehiculeElectrique} dispose de
\texttt{calculerTempsRechargeReel()}, méthode sans équivalent thermique,
illustrant qu'une sous-classe peut encapsuler ses propres règles métier
au-delà du contrat abstrait.}

\section{Programme de test -- \texttt{GestionTransport}}

\enonce{Créer une classe \texttt{GestionTransport} avec \texttt{main()} :
créer 3 véhicules de chaque type, afficher leurs métriques, identifier le
thermique le moins polluant, calculer les autonomies moyennes, créer une
liste unique de tous les véhicules et simuler un trajet de 200 km pour
chacun.}

\screenshot{Screenshot 2026-05-03 162433.png}{Programme de test -- liste polymorphique et affichage des métriques thermiques}
\screenshot{Screenshot 2026-05-03 162506.png}{Programme de test -- métriques des véhicules électriques et comparaison}
\screenshot{Screenshot 2026-05-03 162521.png}{Programme de test -- résultats d'exécution complets}

\begin{lstlisting}[caption={\texttt{GestionTransport.java} -- liste polymorphique et simulation de trajet}]
package TP4;
import java.util.ArrayList;
import java.util.Iterator;
import java.util.List;
import java.util.ListIterator;

public class GestionTransport {
    public static void main(String[] args) {
        // Création des données
        ArrayList<VehiculeThermique> l1 = new ArrayList<>();
        VehiculeThermique t1 = new VehiculeThermique();
        VehiculeThermique t2 = new VehiculeThermique();
        VehiculeThermique t3 = new VehiculeThermique();
        l1.add(t1); l1.add(t2); l1.add(t3);

        ArrayList<VehiculeElectrique> l2 = new ArrayList<>();
        VehiculeElectrique e1 = new VehiculeElectrique();
        VehiculeElectrique e2 = new VehiculeElectrique();
        VehiculeElectrique e3 = new VehiculeElectrique();
        l2.add(e1); l2.add(e2); l2.add(e3);

        // Affichage des métriques pour les véhicules thermiques
        System.out.println("-------Métriques thermiques---------");
        ListIterator<VehiculeThermique> it = l1.listIterator();
        while (it.hasNext()) System.out.println(it.next().toString());

        VehiculeThermique minV = l1.get(0);
        for (VehiculeThermique p : l1)
            if (p.calculerPollution() < minV.calculerPollution()) minV = p;
        System.out.println("Le véhicule thermique le plus écologique : " + minV);

        // Affichage des métriques pour les véhicules électriques
        System.out.println("-------Métriques électriques---------");
        for (VehiculeElectrique p : l2) System.out.println(p);

        double moyE = 0;
        for (VehiculeElectrique p : l2) moyE += p.calculerAutonomie();
        moyE /= l2.size();
        System.out.println("Autonomie moyenne (électriques): " + moyE);

        // Liste polymorphique unique
        double seuil = 10;
        List<Vehicule> liste1 = new ArrayList<>();
        liste1.addAll(l1);
        liste1.addAll(l2);
        Iterator<Vehicule> it3 = liste1.iterator();
        while (it3.hasNext()) {
            Vehicule v3 = it3.next();
            if (v3.calculerEfficacite() > seuil) v3.toString();
        }

        double moyT = 0;
        for (VehiculeThermique p : l1) moyT += p.calculerAutonomie();
        moyT /= l1.size();
        System.out.println("Autonomie moyenne thermiques vs électriques :"
            + "\nthermiques : " + moyT + "\nélectriques: " + moyE);

        // Simulation de trajet de 200 km
        for (Vehicule v : liste1) {
            System.out.println(v.getImmatriculation()
                + (v.calculerAutonomie() >= 200
                    ? " peut effectuer le trajet"
                    : " ne peut pas effectuer le trajet"));
        }
    }
}
\end{lstlisting}

\observation{Le programme stocke des objets de types concrets dans une
\texttt{List<Vehicule>} : c'est le polymorphisme par substitution. L'appel de
\texttt{v.calculerAutonomie()} est résolu dynamiquement à l'exécution
(\textit{late binding}), selon le type réel de l'objet.}

\note{La recherche du véhicule thermique le moins polluant utilise une liste
filtrée de type \texttt{List<VehiculeThermique>}, car \texttt{calculerPollution()}
n'est pas déclarée dans \texttt{Vehicule}.}

\probleme{Depuis la liste de type \texttt{List<Vehicule>}, appeler directement
\texttt{calculerPollution()} provoquait une erreur de compilation car le
compilateur ne connaît que le type déclaré. La solution a été de filtrer d'abord
avec \texttt{instanceof}, puis d'effectuer le cast avant l'appel.}

\conclusionTP{Ce TP approfondit le polymorphisme via les classes abstraites.
La distinction entre méthodes abstraites (contrat imposé) et méthodes concrètes
(comportement hérité) est fondamentale. Le pattern «~liste de type parent,
objets de types enfants~» est l'un des plus utilisés en Java.}

% ════════════════════════════════════════════════════════════════════════════
%  TP 5 : Collections Java
% ════════════════════════════════════════════════════════════════════════════
\tpchapter{5}{Collections Java -- Gestion des Employés}

\objectif{Manipuler \texttt{HashSet}, \texttt{TreeSet},
\texttt{LinkedHashSet}, \texttt{ArrayList}, \texttt{HashMap},
\texttt{Hashtable} ; utiliser \texttt{Iterator} et \texttt{Enumeration} ;
comparer les performances ; convertir entre collections et tableaux.}

\section{Classe \texttt{Employe}}

\enonce{Créer la classe \texttt{Employe} avec les attributs \texttt{id}
(\texttt{int}), \texttt{nom}, \texttt{departement} (\texttt{String}) et
\texttt{salaire} (\texttt{double}). Ajouter : constructeur complet,
getters/setters, \texttt{toString()} formaté, \texttt{equals()} basé sur
\texttt{id}, et implémenter \texttt{Comparable<Employe>} pour trier par
\texttt{id}.}

\screenshot{Screenshot 2026-05-03 162849.png}{Classe \texttt{Employe} -- implémentation de \texttt{Comparable} et \texttt{equals}}

\begin{lstlisting}[caption={\texttt{Employe.java} -- classe métier avec \texttt{Comparable}}]
package TP5;

public class Employe implements Comparable<Employe> {
    private int id;
    private String nom;
    private String departement;
    private double salaire;

    public Employe(int id, String nom, String departement, double salaire) {
        this.id = id;
        this.nom = nom;
        this.departement = departement;
        this.salaire = salaire;
    }

    public int getId()              { return id; }
    public String getNom()          { return nom; }
    public String getDepartement()  { return departement; }
    public double getSlaire()       { return salaire; }

    public void setId(int id)                    { this.id = id; }
    public void setNom(String nom)               { this.nom = nom; }
    public void setepartement(String d)          { this.departement = d; }
    public void setSlaire(double s)              { this.salaire = s; }

    @Override
    public String toString() {
        return "Id : " + id + "; Nom: " + nom
             + "Departement: " + departement
             + "Salaire: " + salaire;
    }

    @Override
    public boolean equals(Object obj) {
        return id == ((Employe) obj).getId();
    }

    @Override
    public int compareTo(Employe autre) {
        return Integer.compare(this.id, autre.id);
    }
}
\end{lstlisting}

\observation{L'implémentation de \texttt{Comparable<Employe>} est obligatoire
pour que \texttt{TreeSet} puisse trier automatiquement les employés. Sans elle,
une \texttt{ClassCastException} serait levée à la première insertion. La
redéfinition d'\texttt{equals()} basée sur \texttt{id} garantit l'unicité dans
les \texttt{Set}.}

\probleme{À la première insertion dans un \texttt{TreeSet<Employe>}, une
\texttt{ClassCastException} était levée car la classe n'implémentait pas
\texttt{Comparable}. L'ajout de \texttt{implements Comparable<Employe>} et de
la méthode \texttt{compareTo()} a résolu le problème.}

\section{Classe principale \texttt{GestionEmployes}}

\enonce{Dans \texttt{GestionEmployes} : créer un \texttt{HashSet} de 5
employés, le copier dans un \texttt{TreeSet}, mesurer les performances
d'affichage, créer un \texttt{LinkedHashSet} fusion des deux, convertir en
tableau, puis répartir dans \texttt{ArrayList} (salaire $> 5\,000$ DH) et
\texttt{Hashtable} (salaire $\leq 5\,000$ DH). Parcourir avec
\texttt{Iterator} et \texttt{Enumeration}. Copier dans \texttt{HashMap} et
comparer les performances \texttt{Hashtable} vs \texttt{HashMap}.}

\screenshot{Screenshot 2026-05-03 162922.png}{\texttt{GestionEmployes} -- initialisation des différentes collections}
\screenshot{Screenshot 2026-05-03 162942.png}{\texttt{GestionEmployes} -- parcours avec \texttt{Iterator} et \texttt{Enumeration}}
\screenshot{Screenshot 2026-05-03 163005.png}{\texttt{GestionEmployes} -- utilisation de \texttt{HashMap} et \texttt{Hashtable}}
\screenshot{Screenshot 2026-05-03 163023.png}{\texttt{GestionEmployes} -- test de performance et résultats}

\begin{lstlisting}[caption={\texttt{GestionEmployes.java} -- panorama complet du framework Collections}]
package TP5;
import java.util.*;

public class GestionEmployes {
    public static void main(String[] arg) {
        // === Collections Set ===
        // HashSet
        HashSet<Employe> ensembleEmployes = new HashSet<>();
        for (int i = 0; i < 5; i++)
            ensembleEmployes.add(new Employe(i, "Etudiant" + i + 1,
                    "D" + i + 1, (i + 1) * 1500));

        // TreeSet (tri naturel via Comparable)
        TreeSet<Employe> ensembleTries = new TreeSet<>(ensembleEmployes);

        // Mesure des performances
        long t0 = System.nanoTime();
        System.out.println(ensembleEmployes);
        long t1 = System.nanoTime();
        System.out.println("temps affichage HashSet : " + (t1 - t0));

        t0 = System.nanoTime();
        System.out.println(ensembleTries);
        t1 = System.nanoTime();
        System.out.println("temps affichage TreeSet : " + (t1 - t0));

        // LinkedHashSet : fusion préservant l'ordre d'insertion
        LinkedHashSet<Employe> ensembleLinked = new LinkedHashSet<>();
        ensembleLinked.addAll(ensembleEmployes);
        ensembleLinked.addAll(ensembleTries);

        // Conversion en tableau
        Employe[] tableauEmployes = new Employe[2];
        tableauEmployes = ensembleLinked.toArray(tableauEmployes);

        // === Collections Map ===
        Hashtable<Integer, String> mapEmployes = new Hashtable<>();
        ArrayList<String> listeNomsEmployes = new ArrayList<>();
        for (int i = 0; i < tableauEmployes.length; i++) {
            if (tableauEmployes[i].getSlaire() > 5000)
                listeNomsEmployes.add(tableauEmployes[i].getNom());
            else
                mapEmployes.put(tableauEmployes[i].getId(),
                                tableauEmployes[i].getNom());
        }

        // Recherche dans la Hashtable
        if (mapEmployes.containsKey(103))      System.out.println("L'id 103 exist");
        if (mapEmployes.containsValue("Mohamed")) System.out.println("Le Nom Mohamed exist");

        // HashMap (copie de la Hashtable)
        HashMap<Integer, String> copieMapEmployes = new HashMap<>(mapEmployes);

        // === Itérateurs ===
        // Iterator (Collection)
        Iterator<String> iterator = listeNomsEmployes.iterator();
        while (iterator.hasNext()) System.out.println(iterator.next());

        // Enumeration (Hashtable)
        Enumeration<String> e = mapEmployes.elements();
        while (e.hasMoreElements()) System.out.println(e.nextElement());

        // entrySet() (HashMap)
        for (Map.Entry<Integer, String> s : copieMapEmployes.entrySet())
            System.out.println("Id : " + s.getKey() + " | Nom : " + s.getValue());

        // === Test de performance Hashtable vs HashMap ===
        t0 = System.nanoTime();
        for (Map.Entry<Integer, String> s : mapEmployes.entrySet())
            System.out.println("Id : " + s.getKey() + " | Nom : " + s.getValue());
        long r0 = System.nanoTime() - t0;

        t0 = System.nanoTime();
        for (Map.Entry<Integer, String> s : copieMapEmployes.entrySet())
            System.out.println("Id : " + s.getKey() + " | Nom : " + s.getValue());
        long r1 = System.nanoTime() - t0;

        System.out.println("différence Hashtable vs HashMap (ns) : " + (r1 - r0));
    }
}
\end{lstlisting}

\observation{Le test avec \texttt{System.nanoTime()} montre que \texttt{HashSet}
est généralement plus rapide que \texttt{TreeSet} (pas de maintien d'ordre).
\texttt{Hashtable} est thread-safe mais plus lente que \texttt{HashMap} en
mono-thread en raison de sa synchronisation inhérente.}

\note{%
Tableau récapitulatif des collections utilisées dans ce TP :\\[6pt]
\begin{tabular}{@{}lllc@{}}
\toprule
\textbf{Collection} & \textbf{Ordre} & \textbf{Doublons} & \textbf{Thread-safe} \\
\midrule
\texttt{HashSet}       & Non garanti    & Non & Non \\
\texttt{TreeSet}       & Trié (naturel) & Non & Non \\
\texttt{LinkedHashSet} & Insertion      & Non & Non \\
\texttt{ArrayList}     & Insertion      & Oui & Non \\
\texttt{HashMap}       & Non garanti    & Clés : Non & Non \\
\texttt{Hashtable}     & Non garanti    & Clés : Non & Oui \\
\bottomrule
\end{tabular}
}

\conclusionTP{Ce TP a permis de comparer concrètement les principales structures
du framework Java Collections. Le choix d'une collection dépend du besoin :
ordre d'insertion, tri naturel, unicité, accès par clé, ou sécurité multi-thread.}

% ════════════════════════════════════════════════════════════════════════════
%  TP 6 : Exceptions personnalisées
% ════════════════════════════════════════════════════════════════════════════
\tpchapter{6}{Exceptions Personnalisées -- Système Bancaire}

\objectif{Comprendre la hiérarchie des exceptions Java, créer une hiérarchie
d'exceptions personnalisées cohérente (\texttt{BanqueException} et cinq
sous-exceptions métier), et les utiliser dans une classe
\texttt{CompteBancaire} robuste.}

\section{Hiérarchie des exceptions bancaires}

\enonce{Créer \texttt{BanqueException} héritant d'\texttt{Exception}, puis
cinq sous-classes : \texttt{SoldeInsuffisantException},
\texttt{CompteInexistantException}, \texttt{MontantInvalideException},
\texttt{DecouvertMaximumDepasseException} et \texttt{CompteBloqueException}.}

\screenshot{Screenshot 2026-05-17 143606.png}{\texttt{BanqueException} -- classe mère des exceptions personnalisées}
\screenshot{Screenshot 2026-05-17 143621.png}{\texttt{SoldeInsuffisantException} et \texttt{CompteInexistantException}}
\screenshot{Screenshot 2026-05-17 143637.png}{\texttt{MontantInvalideException} et \texttt{DecouvertMaximumDepasseException}}
\screenshot{Screenshot 2026-05-17 143653.png}{\texttt{CompteBloqueException}}
\screenshot{Screenshot 2026-05-17 143705.png}{Vue d'ensemble de la hiérarchie d'exceptions}
\screenshot{Screenshot 2026-05-17 143724.png}{Tests de la hiérarchie d'exceptions}

\begin{lstlisting}[caption={Hiérarchie des exceptions personnalisées du système bancaire}]
// === Classe mère ===
package TP6;
public class BanqueException extends Exception {
    public BanqueException(String mssg) { super(mssg); }
}

// === Sous-exceptions métier ===
public class SoldeInsuffisantException        extends BanqueException {
    public SoldeInsuffisantException(String m)        { super(m); } }

public class CompteInexistantException        extends BanqueException {
    public CompteInexistantException(String m)        { super(m); } }

public class MontantInvalideException         extends BanqueException {
    public MontantInvalideException(String m)         { super(m); } }

public class DecouvertMaximumDepasseException extends BanqueException {
    public DecouvertMaximumDepasseException(String m) { super(m); } }

public class CompteBloqueException            extends BanqueException {
    public CompteBloqueException(String m)            { super(m); } }
\end{lstlisting}

\observation{Toutes les exceptions héritent de \texttt{BanqueException}
(elle-même héritière d'\texttt{Exception}). Un bloc
\texttt{catch(BanqueException e)} peut ainsi intercepter toutes les
sous-exceptions, offrant une gestion hiérarchique flexible.}

\note{%
La hiérarchie comprend cinq exceptions métier :
\begin{itemize}[leftmargin=1.5cm, itemsep=2pt]
  \item \texttt{SoldeInsuffisantException} -- retrait supérieur au solde disponible
  \item \texttt{CompteInexistantException} -- compte introuvable dans le registre
  \item \texttt{MontantInvalideException} -- montant négatif ou nul
  \item \texttt{DecouvertMaximumDepasseException} -- découvert autorisé dépassé
  \item \texttt{CompteBloqueException} -- opération sur un compte bloqué
\end{itemize}
}

\section{Classe \texttt{CompteBancaire}}

\enonce{Implémenter \texttt{CompteBancaire} avec les attributs
\texttt{numeroCompte}, \texttt{titulaire}, \texttt{solde},
\texttt{decouvertAutorise}, \texttt{bloque}, \texttt{dateCreation} et
\texttt{historiqueOperations} (\texttt{ArrayList<String>}). Créer les
méthodes : \texttt{deposer()}, \texttt{retirer()},
\texttt{transferer()}, \texttt{afficherHistorique()},
\texttt{ajouterCompte()}, \texttt{getCompte()},
\texttt{effectuerVirement()} et \texttt{afficherTousLesComptes()}.}

\screenshot{Screenshot 2026-05-17 142718.png}{\texttt{CompteBancaire} -- attributs, constructeur et registre statique des comptes}
\screenshot{Screenshot 2026-05-17 142734.png}{\texttt{CompteBancaire} -- méthode \texttt{deposer()} avec validation}
\screenshot{Screenshot 2026-05-17 142756.png}{\texttt{CompteBancaire} -- méthode \texttt{retirer()} avec gestion complète des exceptions}
\screenshot{Screenshot 2026-05-17 142820.png}{\texttt{CompteBancaire} -- méthode \texttt{effectuerVirement()} et journalisation}
\screenshot{Screenshot 2026-05-17 142837.png}{\texttt{CompteBancaire} -- méthodes utilitaires et \texttt{toString()}}

\begin{lstlisting}[caption={\texttt{CompteBancaire.java} -- attributs, registre statique et méthodes (extrait)}]
package TP6;
import java.util.*;

public class CompteBancaire {
    private String numeroCompte;
    private String titulaire;
    private double solde;
    private double decouvertAutorise;
    private boolean bloque;
    private Date dateCreation;
    private ArrayList<String> historiqueOperations;
    private static List<CompteBancaire> listeTousLesComptes = new ArrayList<>();

    public CompteBancaire(String numeroCompte, String titulaire, double solde,
                          double decouvertAutorise, boolean bloque, Date dateCreation) {
        try {
            for (CompteBancaire c : getTousLesComptes())
                if (c.getNumeroCompte().equals(numeroCompte))
                    throw new BanqueException("Le compte du numero "
                            + numeroCompte + " est deja exist");
            this.numeroCompte = numeroCompte;
        } catch (BanqueException e) {
            System.out.println(e.getMessage());
            return;
        }
        this.titulaire = titulaire;
        this.solde = solde;
        this.decouvertAutorise = decouvertAutorise;
        this.bloque = bloque;
        this.dateCreation = dateCreation;
        this.historiqueOperations = new ArrayList<>();
        listeTousLesComptes.add(this);
    }

    // --- Getters/Setters omis pour la concision ---

    public void deposer(double nbr) {
        try {
            if (nbr < 0)        throw new MontantInvalideException("le montant doit etre positif");
            else if (getBloque()) throw new CompteBloqueException("le compte est bloqué");
            setSolde(getSolde() + nbr);
            setHistoriqueOperations("Deposer: " + nbr);
        } catch (MontantInvalideException | CompteBloqueException e) {
            System.out.println(e.getMessage());
        }
    }

    public void retirer(double nbr) {
        try {
            if (nbr < 0)        throw new MontantInvalideException("le montant doit etre positif");
            else if (getBloque()) throw new CompteBloqueException("le compte est bloqué");
            else if (getSolde() - nbr < getDecouvertAutorise())
                throw new DecouvertMaximumDepasseException("Insufisant");
            setSolde(getSolde() - nbr);
            setHistoriqueOperations("Retirer: " + nbr);
        } catch (MontantInvalideException | CompteBloqueException
               | DecouvertMaximumDepasseException e) {
            System.out.println(e.getMessage());
        }
    }

    public static void effectuerVirement(String numeroSource,
                                         String numeroDestinataire, double montant) {
        CompteBancaire s = getCompte(numeroSource);
        CompteBancaire d = getCompte(numeroDestinataire);
        if (s != null && d != null) {
            s.retirer(montant);
            d.deposer(montant);
        }
    }
}
\end{lstlisting}

\observation{Le registre statique \texttt{listeTousLesComptes} centralise tous
les comptes créés, permettant la recherche par numéro depuis n'importe quel
point du programme. La méthode \texttt{effectuerVirement()} délègue à
\texttt{retirer()} et \texttt{deposer()}, réutilisant la gestion des exceptions
déjà en place -- c'est le principe DRY (\textit{Don't Repeat Yourself}).}

\erreurFrequente{Lancer une exception \textit{checked} (héritant d'
\texttt{Exception}) sans la déclarer avec \texttt{throws} dans la signature,
ni la capturer avec \texttt{catch}, provoque une erreur de compilation.
Elle doit être soit gérée localement, soit propagée explicitement.}

\probleme{Lors de l'implémentation d'\texttt{effectuerVirement()}, l'oubli de
la clause \texttt{throws BanqueException} dans la signature empêchait la
compilation. En propageant correctement les exceptions via \texttt{throws},
le compilateur a accepté le code et la gestion s'est faite au niveau du
programme principal.}

\section{Programme principal -- \texttt{GestionCompteBancaire}}

\enonce{Dans \texttt{main()} : ajouter 3 comptes, tenter d'ajouter un
numéro existant, faire des dépôts/retraits, effectuer un virement, tester
tous les cas d'erreur (solde insuffisant, montant négatif, compte inexistant,
compte bloqué).}

\screenshot{Screenshot 2026-05-17 144059.png}{Programme principal -- tests de toutes les opérations bancaires et gestion des exceptions}

\begin{lstlisting}[caption={\texttt{GestionCompteBancaire.java} -- scénarios de test couvrant tous les cas d'erreur}]
package TP6;
import java.util.Date;

public class GestionCompteBancaire {
    public static void main(String[] args) {
        // 3 comptes valides
        CompteBancaire c1 = new CompteBancaire("C1", "yahya", 100,  50, false, new Date());
        CompteBancaire c2 = new CompteBancaire("C2", "adam",  1000, 50, false, new Date());
        CompteBancaire c3 = new CompteBancaire("C3", "ayoub", 500,  50, false, new Date());

        // Doublon : numéro existant -> CompteInexistantException levée
        CompteBancaire c4 = new CompteBancaire("C3", "ayoub", 500, 50, false, new Date());

        c1.deposer(200);                       // dépôt valide
        c2.retirer(900);                       // retrait valide
        CompteBancaire.effectuerVirement("C3", "C1", 200);  // virement OK
        c2.retirer(60);                        // découvert max
        c3.deposer(-1);                        // montant négatif
        CompteBancaire.effectuerVirement("C3", "C7", 10);   // compte inexistant

        // Compte bloqué
        CompteBancaire c5 = new CompteBancaire("C5", "ayoub", 500, 50, true, new Date());
        c5.deposer(100);
    }
}
\end{lstlisting}

\conclusionTP{Ce TP démontre la valeur des exceptions personnalisées pour
produire des messages d'erreur métier précis. La séparation entre logique de
validation (\texttt{CompteBancaire}) et programme principal simplifie la
maintenance et améliore la robustesse du système bancaire.}

% ════════════════════════════════════════════════════════════════════════════
%  TP 7 : Threads et synchronisation
% ════════════════════════════════════════════════════════════════════════════
\tpchapter{7}{Threads et Synchronisation -- Distributeur Automatique de Billets}

\objectif{Maîtriser la création et la gestion des threads Java, comprendre
les conditions de course (\textit{race conditions}) en accès concurrent, et
utiliser \texttt{synchronized}, \texttt{wait()} et \texttt{notifyAll()} pour
protéger les ressources partagées.}

\section{Partie 1 -- Sans synchronisation : mise en évidence du problème}

\enonce{Créer une classe \texttt{Distributeur} (fonds = 5\,000 DH) avec
une méthode \texttt{retirer(montant, nomClient)} incluant un
\texttt{Thread.sleep(2000)} pour simuler le traitement. Créer la classe
\texttt{Client} (\texttt{Runnable}) effectuant $n$ retraits en boucle.
Dans \texttt{Main} : deux clients \texttt{Ahmed} et \texttt{Amal},
chacun tente 3 retraits de 3\,000 DH. Observer les incohérences.}

\screenshot{Screenshot 2026-05-17 145750.png}{\texttt{Distributeur} non synchronisé -- vérification et retrait sans protection}
\screenshot{Screenshot 2026-05-17 145807.png}{Classe \texttt{Client} (\texttt{Runnable}) et classe \texttt{Main} de lancement}
\screenshot{Screenshot 2026-05-17 145828.png}{Résultat d'exécution -- condition de course : solde négatif observé}

\begin{lstlisting}[caption={\texttt{Distributeur.java} -- version non synchronisée (illustre la condition de course)}]
package TP7;

public class Distributeur {
    private double fonds = 5000;
    public double getFonds()              { return fonds; }
    public void   setFonds(double fonds)  { this.fonds = fonds; }

    public void retirer(double montant, String nomClient) {
        System.out.println(nomClient + " vérifie le solde : " + getFonds());
        if (montant <= getFonds()) {
            System.out.println("Solde suffisant.");
            setFonds(getFonds() - montant);
            try { Thread.sleep(2000); }
            catch (InterruptedException e) { throw new RuntimeException(e); }
            System.out.println("Le nouveau solde : " + getFonds());
        } else
            System.out.println("Solde insuffisant.");
    }
}
\end{lstlisting}

\begin{lstlisting}[caption={\texttt{Client.java} -- thread client effectuant n retraits}]
package TP7;

class Client implements Runnable {
    Distributeur distributeur;
    String nom;
    double montantRetrait;
    int nombreOperations;

    public Client(Distributeur d, String n, double m, int nbr) {
        distributeur = d; nom = n; montantRetrait = m; nombreOperations = nbr;
    }

    @Override
    public void run() {
        int i = nombreOperations;
        while (i > 0) {
            distributeur.retirer(montantRetrait, nom);
            try { Thread.sleep(1000); }
            catch (InterruptedException e) { throw new RuntimeException(e); }
            i--;
        }
    }
}
\end{lstlisting}

\begin{lstlisting}[caption={\texttt{Test.java} -- lancement de deux threads concurrents}]
package TP7;

public class Test {
    public static void main(String[] args) {
        Distributeur d = new Distributeur();
        Client c1 = new Client(d, "Ahmed", 3000, 3);
        Client c2 = new Client(d, "Amal",  3000, 3);
        Thread t1 = new Thread(c1);
        Thread t2 = new Thread(c2);
        t1.start(); t2.start();
    }
}
\end{lstlisting}

\observation{Sans synchronisation, une \textbf{condition de course}
(\textit{race condition}) survient : les deux threads lisent simultanément
\texttt{fonds = 5\,000}, croient tous les deux avoir assez, et retirent
chacun 3\,000 DH, portant le solde à $-1\,000$ DH. Le
\texttt{Thread.sleep(2000)} amplifie intentionnellement ce problème en
simulant un délai entre vérification et débit.}

\erreurFrequente{Lire une variable partagée puis la modifier en deux
instructions distinctes sans section critique crée systématiquement des
incohérences en environnement multi-thread. Ce n'est pas une opération
atomique.}

\probleme{En exécutant la version non synchronisée plusieurs fois, le solde
final variait à chaque exécution (tantôt $-1\,000$ DH, tantôt $+2\,000$ DH,
tantôt $+5\,000$ DH). Ce comportement non déterministe a permis de comprendre
concrètement ce qu'est une condition de course.}

\section{Partie 2 -- Avec synchronisation : solution correcte}

\enonce{Créer \texttt{DistributeurSecurise} avec une méthode
\texttt{synchronized retirer()} utilisant \texttt{while + wait()} et une
méthode \texttt{synchronized deposer()} avec \texttt{notifyAll()}. Créer
\texttt{ClientOperation} (\texttt{Runnable}). Dans \texttt{MainSecurise} :
Ahmed (3 retraits de 3\,000 DH), Amal (2 dépôts de 2\,000 DH + 1 retrait
de 3\,000 DH), et un thread démon afficheur toutes les 3 secondes.}

\screenshot{Screenshot 2026-05-17 150146.png}{\texttt{DistributeurSecurise} -- méthodes \texttt{retirer()} et \texttt{deposer()} synchronisées}
\screenshot{Screenshot 2026-05-17 150203.png}{Classe \texttt{ClientOperation} et thread démon afficheur}
\screenshot{Screenshot 2026-05-17 150226.png}{Résultat d'exécution -- solde cohérent garanti par la synchronisation}

\begin{lstlisting}[caption={\texttt{DistributeurSecurise.java} -- version thread-safe avec \texttt{synchronized + wait/notifyAll}}]
package TP7;

public class DistributeurSecurise {
    private double fonds = 5000;

    public synchronized double getFonds()              { return fonds; }
    public synchronized void   setFonds(double fonds)  { this.fonds = fonds; }

    public synchronized void retirer(double montant, String nomClient) {
        System.out.println(nomClient + " tente de retirer : " + montant);
        while (montant > getFonds()) {                 // boucle while OBLIGATOIRE
            System.out.println("Solde insuffisant, " + nomClient + " attend...");
            try { wait(); }
            catch (InterruptedException e) { throw new RuntimeException(e); }
        }
        try { Thread.sleep(2000); }
        catch (InterruptedException e) { throw new RuntimeException(e); }
        setFonds(getFonds() - montant);
        System.out.println("Le nouveau solde : " + getFonds());
    }

    public synchronized void deposer(double montant, String nomClient) {
        System.out.println(nomClient + " dépose : " + montant);
        setFonds(getFonds() + montant);
        System.out.println("Le nouveau solde : " + getFonds());
        notifyAll();                                   // réveille les threads en attente
    }
}
\end{lstlisting}

\begin{lstlisting}[caption={\texttt{ClientOperation.java} -- thread polyvalent (dépôt ou retrait)}]
package TP7;

public class ClientOperation implements Runnable {
    DistributeurSecurise distributeur;
    String nom;
    boolean estDepot;
    double montant;
    int nombreOperations;

    public ClientOperation(DistributeurSecurise d, String n, boolean i, double m, int nbr) {
        distributeur = d; nom = n; estDepot = i; montant = m; nombreOperations = nbr;
    }

    @Override
    public void run() {
        int i = nombreOperations;
        while (i > 0) {
            if (estDepot) distributeur.deposer(montant, nom);
            else          distributeur.retirer(montant, nom);
            try { Thread.sleep(1000); }
            catch (InterruptedException e) { throw new RuntimeException(e); }
            i--;
        }
    }
}
\end{lstlisting}

\begin{lstlisting}[caption={\texttt{MainSecurise.java} -- orchestration des threads et thread démon afficheur}]
package TP7;

public class MainSecurise {
    public static void main(String[] args) {
        DistributeurSecurise d = new DistributeurSecurise();

        ClientOperation c1 = new ClientOperation(d, "Ahmed", false, 3000, 3);
        ClientOperation c2 = new ClientOperation(d, "Amal",  true,  2000, 2);
        ClientOperation c3 = new ClientOperation(d, "Amal",  false, 3000, 1);

        Thread t1 = new Thread(c1);
        Thread t2 = new Thread(c2);
        Thread t3 = new Thread(c3);

        Runnable afficheur = () -> {
            while (true) {
                System.out.println("Solde actuel : " + d.getFonds());
                try { Thread.sleep(3000); }
                catch (InterruptedException e) { return; }
            }
        };
        Thread c = new Thread(afficheur, "afficheur");
        c.setDaemon(true);                  // se termine avec le programme
        c.start();

        t1.start(); t2.start(); t3.start();
    }
}
\end{lstlisting}

\observation{L'utilisation de \texttt{while (fonds < montant) \{ wait(); \}}
est préférable à \texttt{if} car un thread peut être réveillé par
\texttt{notifyAll()} sans que la condition soit réellement satisfaite
(\textit{spurious wakeup}). La boucle \texttt{while} force une
re-vérification systématique.}

\note{Le thread \texttt{afficheur} est déclaré \texttt{daemon} : il se termine
automatiquement quand tous les threads non-daemon ont terminé, évitant tout
blocage indéfini du programme.}

\probleme{Dans la première implémentation synchronisée, \texttt{if} était
utilisé à la place de \texttt{while} dans la condition d'attente. Le thread
se réveillait parfois alors que les fonds étaient toujours insuffisants,
provoquant des retraits incorrects. Remplacer \texttt{if} par \texttt{while}
a résolu le problème en forçant une re-vérification à chaque réveil.}

\conclusionTP{Ce TP illustre deux réalités fondamentales de la concurrence :
(1)~l'absence de synchronisation produit des résultats imprévisibles ;
(2)~\texttt{synchronized} + \texttt{wait/notifyAll} constitue le mécanisme
de base Java pour protéger les ressources partagées et coordonner les threads.}

% ════════════════════════════════════════════════════════════════════════════
%  TP 9 : Gestion des fichiers et flux
% ════════════════════════════════════════════════════════════════════════════
\tpchapter{9}{Gestion des Fichiers et Flux}

\objectif{Maîtriser les trois niveaux d'accès aux fichiers en Java :
(1)~fichiers texte avec \texttt{BufferedReader/PrintWriter},
(2)~fichiers binaires typés avec \texttt{DataInputStream/DataOutputStream},
(3)~persistance d'objets avec la sérialisation (\texttt{ObjectOutputStream}).}

\section{Exercice 1 -- Catalogue produits (fichier texte)}

\enonce{Créer un fichier \texttt{catalogue.txt}, y écrire une liste de
produits (nom, prix unitaire), puis le lire et afficher son contenu sur la
console.}

\screenshot{Screenshot 2026-05-17 150625.png}{Code -- Écriture et lecture d'un catalogue texte avec \texttt{try-with-resources}}

\begin{lstlisting}[caption={Catalogue produits -- écriture/lecture d'un fichier texte}]
// Ecriture
try (BufferedWriter bw = new BufferedWriter(new FileWriter("./catalogue.txt"))) {
    bw.write("nom,prix unitaire"); bw.newLine();
    bw.write("Stylo,5.00");        bw.newLine();
    bw.write("Cahier,15.00");
} catch (IOException e) { throw new RuntimeException(e); }

// Lecture
try (BufferedReader br = new BufferedReader(new FileReader("./catalogue.txt"))) {
    String line;
    while ((line = br.readLine()) != null)
        System.out.println(line);
} catch (IOException e) { throw new RuntimeException(e); }
\end{lstlisting}

\observation{L'instruction \texttt{try-with-resources} (Java 7+) garantit la
fermeture automatique des flux même en cas d'exception, éliminant les fuites
de ressources. \texttt{PrintWriter} wrappé autour de \texttt{FileWriter} ajoute
des méthodes de haut niveau comme \texttt{println()}.}

\probleme{Sans fermeture explicite des flux, le fichier n'était pas entièrement
écrit sur le disque (données encore en mémoire tampon). En adoptant
\texttt{try-with-resources}, la fermeture et le vidage du tampon sont
automatiquement garantis à la sortie du bloc.}

\section{Exercice 2 -- Analyse CSV de commandes}

\enonce{Lire un fichier \texttt{commandes.csv} contenant les colonnes
\texttt{id\_commande, nom\_client, produit, quantite, prix\_unitaire},
calculer le montant total de chaque commande (quantité $\times$ prix
unitaire) et afficher les données dans un tableau formaté.}

\screenshot{Screenshot 2026-05-17 150649.png}{Code -- Parsing CSV ligne par ligne et calcul des totaux par commande}

\begin{lstlisting}[caption={\texttt{Test.java} (TP9) -- analyse d'un fichier CSV de commandes}]
package TP9;
import java.io.*;
import java.util.Arrays;
import java.util.List;

public class Test {
    public static void main(String[] args) {
        try (BufferedReader cmd = new BufferedReader(
                new FileReader("./src/TP9/commande.csv"))) {
            String line;
            List<String> l;
            float total;

            line = cmd.readLine();   // saute l'en-tête
            if (line != null) {
                System.out.println("id_commande, montant total");
                while ((line = cmd.readLine()) != null) {
                    l = Arrays.stream(line.split(",")).toList();
                    total = Float.parseFloat(l.get(3))
                          * Float.parseFloat(l.get(4));
                    System.out.println(l.get(0) + ": " + total);
                }
            }
        } catch (IOException e) { throw new RuntimeException(e); }
    }
}
\end{lstlisting}

\observation{\texttt{String.split(",")} suppose un séparateur virgule sans
espace. L'appel à \texttt{.trim()} sur chaque colonne permet de gérer les
espaces parasites. Pour des CSV complexes (guillemets, virgules dans les
valeurs), une bibliothèque dédiée comme \textit{Apache Commons CSV} est
préférable.}

\section{Exercice 3 -- Copie de fichier texte (\texttt{CopieFichier})}

\enonce{Programmer la classe \texttt{CopieFichier} qui copie
\texttt{facture\_src.txt} vers \texttt{facture\_dest.txt} selon deux
méthodes : (1)~caractère par caractère avec
\texttt{FileReader/FileWriter}, (2)~ligne par ligne avec
\texttt{BufferedReader/PrintWriter}. Afficher le contenu source et le nombre
de lignes copiées.}

\screenshot{Screenshot 2026-05-17 150908.png}{\texttt{CopieFichier} -- Méthode 1 : copie caractère par caractère (\texttt{FileReader/FileWriter})}
\screenshot{Screenshot 2026-05-17 150944.png}{\texttt{CopieFichier} -- Méthode 2 : copie ligne par ligne (\texttt{BufferedReader/PrintWriter})}
\screenshot{Screenshot 2026-05-17 150954.png}{Programme de test -- comparaison des deux méthodes et mesure du temps d'exécution}

\begin{lstlisting}[caption={\texttt{CopieFichier.java} -- deux méthodes de copie comparées}]
package TP9;
import java.io.*;

public class CopieFichier {

    // Méthode 1 : copie LIGNE PAR LIGNE (bufferisée -- rapide)
    public static int CopierLigne(String src, String dest) {
        int nbligne = 0;
        try (BufferedReader in  = new BufferedReader(new FileReader(src));
             BufferedWriter out = new BufferedWriter(new FileWriter(dest))) {
            String line;
            while ((line = in.readLine()) != null) {
                out.write(line);
                out.newLine();
                nbligne++;
            }
        } catch (IOException e) { throw new RuntimeException(e); }
        return nbligne;
    }

    // Méthode 2 : copie CARACTÈRE PAR CARACTÈRE (non bufferisée -- lente)
    public static int CopierCar(String src, String dest) {
        int nbligne = 0;
        try (FileReader in  = new FileReader(src);
             FileWriter out = new FileWriter(dest)) {
            int c;
            while ((c = in.read()) != -1) {
                out.write(c);
                nbligne++;
            }
        } catch (IOException e) { throw new RuntimeException(e); }
        return nbligne;
    }

    public static void afficherContenue(String src) {
        System.out.println("=== Contenu de : " + src + " ===");
        try (BufferedReader rd = new BufferedReader(new FileReader(src))) {
            String line;
            while ((line = rd.readLine()) != null)
                System.out.println(line);
        } catch (IOException e) { throw new RuntimeException(e); }
    }

    public static void main(String[] args) {
        afficherContenue("./src/TP9/facture_src.txt");
        int n1 = CopierCar  ("./src/TP9/facture_src.txt", "./src/TP9/facture_dest.txt");
        System.out.println("[M1] Lignes copiées : " + n1);
        int n2 = CopierLigne("./src/TP9/facture_src.txt", "./src/TP9/facture_dest.txt");
        System.out.println("[M2] Lignes copiées : " + n2);
    }
}
\end{lstlisting}

\observation{La copie caractère par caractère est nettement plus lente car elle
effectue un appel système pour chaque caractère lu. La copie ligne par ligne,
bufferisée, regroupe les lectures en blocs, ce qui apporte un gain de
performance considérable sur de gros fichiers.}

\note{Sur Windows, les fins de lignes sont \texttt{"\textbackslash r\textbackslash n"}.
La détection de \texttt{'\textbackslash n'} seul peut donner un comptage inexact.
\texttt{BufferedReader.readLine()} gère transparentement les deux formats.}

\probleme{La copie caractère par caractère s'avérait environ 50 fois plus lente
sur un fichier de 100 Ko. Ce résultat illustre concrètement l'importance du
buffering : chaque \texttt{read()} sur un \texttt{FileReader} non bufferisé
déclenche un appel système direct au disque.}

\section{Exercice 4 -- Fichier binaire : \texttt{MultiplesReel}}

\enonce{Créer la classe \texttt{MultiplesReel} (attribut \texttt{base}
de type \texttt{double}). Méthode \texttt{ecrireMultiples(nomFichier, n)} :
écrire les $n$ multiples de \texttt{base} (chaque couple : entier
coefficient + double résultat). Méthode \texttt{lireEtAfficher(nomFichier)} :
déduire la valeur de \texttt{base} et le nombre de multiples stockés.}

\screenshot{Screenshot 2026-05-17 151447.png}{\texttt{MultiplesReel} -- écriture des couples (coefficient, multiple) avec \texttt{DataOutputStream}}
\screenshot{Screenshot 2026-05-17 151501.png}{\texttt{MultiplesReel} -- lecture avec \texttt{DataInputStream} et déduction de \texttt{base}}

\begin{lstlisting}[caption={\texttt{MultiplesReel.java} -- écriture et lecture binaire typée}]
package TP9;
import java.io.*;

public class MultiplesReel {
    private double base;

    public MultiplesReel(double d) { base = d; }

    public void ecrireMultiples(String nomFichier, int n) {
        try (DataOutputStream ou = new DataOutputStream(
                new FileOutputStream(nomFichier))) {
            for (int i = 1; i <= n; i++) {
                double result = i * base;
                ou.write(i);              // coefficient
                ou.writeDouble(result);   // multiple
            }
        } catch (IOException e) { /* gestion silencieuse */ }
    }

    public void lireEtAfficher(String nomFichier) {
        try (DataInputStream in = new DataInputStream(
                new FileInputStream(nomFichier))) {
            int cout = 0;
            double baseD = 0;
            int coeff      = in.readInt();
            double result  = in.readFloat();
            baseD = result / coeff;
            cout++;
            while (true) {
                try { in.readInt(); in.readDouble(); cout++; }
                catch (EOFException e) { break; }
            }
            System.out.println("La valeur du réel base : " + baseD);
            System.out.println("Le nombre de multiples stockés : " + cout);
        } catch (IOException e) {
            System.err.println("Erreur lecture : " + e.getMessage());
        }
    }

    public static void main(String[] arg) {
        MultiplesReel m = new MultiplesReel(2.2);
        m.ecrireMultiples("./src/TP9/facture_dest.txt", 10);
        m.lireEtAfficher("./src/TP9/facture_dest.txt");
    }
}
\end{lstlisting}

\observation{Les fichiers binaires écrits avec \texttt{DataOutputStream}
stockent les données dans un format compact non lisible par un éditeur de
texte, mais portable entre toutes les plateformes Java (format \textit{big-endian}).
Chaque \texttt{int} occupe 4 octets, chaque \texttt{double} 8 octets.}

\erreurFrequente{L'ordre de lecture avec \texttt{DataInputStream} doit être
\emph{exactement} identique à l'ordre d'écriture. Lire un \texttt{double} là
où un \texttt{int} a été écrit produit des valeurs corrompues sans lever
d'exception.}

\probleme{Lors de la lecture, une tentative de lire un \texttt{double} avant
l'\texttt{int} de coefficient produisait des valeurs aberrantes. L'ordre de
lecture doit correspondre exactement à l'ordre d'écriture -- contrainte forte
des flux binaires typés en Java.}

\section{Exercice 5 -- Sérialisation d'objets : \texttt{Client}}

\enonce{Créer la classe \texttt{Client} implémentant \texttt{Serializable}
(attributs : \texttt{nom}, \texttt{prenom}, \texttt{adresse},
\texttt{idClient}). Sauvegarder 3 clients dans \texttt{clients.ser} via
\texttt{ObjectOutputStream}. Recharger depuis \texttt{clients.ser} et
extraire les noms complets (\texttt{prenom + " " + nom}) dans un tableau
de chaînes.}

\screenshot{Screenshot 2026-05-17 151850.png}{Classe \texttt{Client} -- implémentation de \texttt{Serializable} et \texttt{serialVersionUID}}
\screenshot{Screenshot 2026-05-17 151905.png}{Programme de sérialisation -- écriture dans \texttt{clients.ser} et relecture des noms complets}

\begin{lstlisting}[caption={\texttt{Client.java} (TP9) -- objet sérialisable}]
package TP9;
import java.io.Serializable;

public class Client implements Serializable {
    private String nom;
    private String prenom;
    private String adress;
    private int idClient;

    public String getNom()    { return nom; }
    public String getPrenom() { return prenom; }
    public String getAdress() { return adress; }

    public Client(int id, String n, String p, String a) {
        idClient = id; nom = n; prenom = p; adress = a;
    }
}
\end{lstlisting}

\begin{lstlisting}[caption={\texttt{SerializationDemo.java} -- écriture/lecture d'objets via \texttt{ObjectOutputStream}}]
package TP9;
import java.io.*;

public class SerializationDemo {
    public static void main(String[] args)
            throws IOException, ClassNotFoundException {
        Client[] clients = {
            new Client(1, "Ali",    "Benali",  "Casablanca"),
            new Client(2, "Fatima", "El Amri", "Rabat"),
            new Client(3, "Omar",   "Haddi",   "Khouribga")
        };

        // Sérialisation
        try (ObjectOutputStream oos = new ObjectOutputStream(
                new FileOutputStream("./src/TP9/client.ser"))) {
            for (Client c : clients) oos.writeObject(c);
        }

        // Désérialisation
        String[] noms = new String[clients.length];
        try (ObjectInputStream ois = new ObjectInputStream(
                new FileInputStream("./src/TP9/client.ser"))) {
            for (int i = 0; i < noms.length; i++) {
                Client c = (Client) ois.readObject();
                noms[i] = c.getPrenom() + " " + c.getNom();
            }
        }

        System.out.println("Noms complets : ");
        for (String n : noms) System.out.println(" " + n);
    }
}
\end{lstlisting}

\observation{La constante \texttt{serialVersionUID = 1L} est une bonne pratique
explicite : elle contrôle la compatibilité entre versions de la classe. Sans UID
déclaré, Java en génère un automatiquement qui peut changer à chaque
compilation, rendant les anciens fichiers \texttt{.ser} illisibles.}

\note{La sérialisation Java est pratique mais présente des limites : elle est
spécifique à Java, couplée à la structure interne de la classe, et peut poser
des problèmes de sécurité. Pour les échanges inter-systèmes, JSON ou XML sont
préférables.}

\probleme{Lors des tests de désérialisation, une
\texttt{InvalidClassException} était levée car la classe \texttt{Client} avait
été modifiée entre l'écriture et la lecture du fichier. Java avait généré un
UID différent. L'ajout explicite de
\texttt{private static final long serialVersionUID = 1L} a stabilisé la
compatibilité.}

\conclusionTP{Ce TP a couvert les trois niveaux d'accès aux fichiers en Java :
texte structuré (\texttt{BufferedReader/PrintWriter}), données binaires typées
(\texttt{DataInputStream/DataOutputStream}), et persistance d'objets complets
(sérialisation). Chaque approche a son domaine d'application selon la nature
et l'usage des données.}

% ════════════════════════════════════════════════════════════════════════════
%  TP 10 : Interfaces graphiques et base de données
% ════════════════════════════════════════════════════════════════════════════
\tpchapter{10}{Interfaces Graphiques et Base de Données -- Swing \& JDBC}

\objectif{Maîtriser la création d'interfaces graphiques avec les composants
Java Swing (\texttt{JFrame}, \texttt{JLabel}, \texttt{JTextField},
\texttt{JComboBox}, \texttt{JButton}), et la connexion à une base de données
relationnelle via JDBC pour effectuer des opérations CRUD (\textit{Create,
Read, Update, Delete}).}

\section{Contexte et spécifications}

\enonce{Créer une interface graphique permettant d'enregistrer les
informations des membres d'un club dans une base de données. Les
informations sont stockées dans une table \texttt{T\_Membre}. En utilisant
les composants Swing, réaliser le formulaire avec les champs :
\texttt{IdMembre}, \texttt{Prenom}, \texttt{Nom}, \texttt{Adresse},
\texttt{Date de Naissance}, \texttt{Lieu de Naissance}, \texttt{Sexe}
(liste déroulante : M/F) et \texttt{Profession}. Prévoir trois boutons :
\textbf{Ajouter}, \textbf{Modifier} et \textbf{Quitter}.}

\screenshot{Screenshot 2026-05-TP10-interface.png}{Interface graphique Swing -- Formulaire d'inscription des membres du club}

\observation{La classe \texttt{JFrame} constitue la fenêtre principale de
l'application. Le gestionnaire de disposition (\textit{layout manager})
\texttt{GridBagLayout} ou \texttt{BorderLayout} + \texttt{GridLayout}
permet d'aligner proprement les étiquettes (\texttt{JLabel}) et les
champs de saisie (\texttt{JTextField}). Le composant \texttt{JComboBox}
gère la liste déroulante pour le champ \texttt{Sexe}.}

\section{Structure de la table \texttt{T\_Membre}}

\enonce{Concevoir la table SQL \texttt{T\_Membre} correspondant aux
champs du formulaire, avec \texttt{IdMembre} comme clé primaire, puis
implémenter la connexion JDBC et les méthodes d'insertion et de
modification.}

\screenshot{Screenshot 2026-05-TP10-table.png}{Structure SQL de la table \texttt{T\_Membre} et connexion JDBC}

\observation{La connexion JDBC suit le patron classique en quatre étapes :
(1)~charger le pilote (\texttt{Class.forName(...)}),
(2)~ouvrir la connexion (\texttt{DriverManager.getConnection(...)}),
(3)~créer et exécuter une requête (\texttt{PreparedStatement}),
(4)~fermer les ressources dans un bloc \texttt{finally} ou
\texttt{try-with-resources}. L'utilisation de \texttt{PreparedStatement}
au lieu de \texttt{Statement} protège contre les injections SQL.}

\note{La liste déroulante \texttt{JComboBox<String>} est initialisée avec
les valeurs \texttt{"M"} et \texttt{"F"}. La valeur sélectionnée est
récupérée via \texttt{getSelectedItem().toString()} lors du clic sur
\textbf{Ajouter} ou \textbf{Modifier}.}

\section{Implémentation du formulaire Swing}

\screenshot{Screenshot 2026-05-TP10-code1.png}{Classe principale -- Déclaration de la \texttt{JFrame} et des composants Swing}
\screenshot{Screenshot 2026-05-TP10-code2.png}{Méthode \texttt{actionPerformed} -- Bouton Ajouter et connexion JDBC}
\screenshot{Screenshot 2026-05-TP10-code3.png}{Bouton Modifier -- Requête UPDATE et bouton Quitter}

\begin{lstlisting}[caption={\texttt{Test.java} -- vérification de la connexion MySQL}]
package TP10;
import java.sql.Connection;
import java.sql.DriverManager;
import java.sql.SQLException;

public class Test {
    public static void main(String[] args) {
        String url  = "jdbc:mysql://localhost:3306/t-membre";
        String user = "root";
        String pwd  = "yahya123";
        try (Connection conn = DriverManager.getConnection(url, user, pwd)) {
            System.out.println("Connexion MySQL réussie !");
        } catch (SQLException e) {
            System.out.println("Erreur : " + e.getMessage());
        }
    }
}
\end{lstlisting}

\begin{lstlisting}[caption={\texttt{fenetrePrincipale.java} -- construction de l'interface Swing (extrait)}]
package TP10;
import javax.swing.*;
import javax.swing.border.EmptyBorder;
import java.awt.*;
import java.sql.*;
import java.awt.event.ActionEvent;
import java.awt.event.ActionListener;

public class fenetrePrincipale extends JFrame {

    JPanel titlePanel, formPanel, buttonPanel;
    JLabel titre;
    JLabel lid, lprenom, lnom, ladresse, ldateN, llieuN, lsexe, lprofession;
    JTextField id, prenom, nom, adresse, dateN, lieuN, profession;
    JComboBox<String> sexe;
    JButton bAjouter, bModifier, bQuitter;
    Connection conn;

    public fenetrePrincipale() {
        // === Connexion à MySQL ===
        try {
            String url = "jdbc:mysql://localhost:3306/t-membre"
                       + "?useSSL=false&allowPublicKeyRetrieval=true";
            conn = DriverManager.getConnection(url, "root", "yahya123");
            System.out.println("Connexion réussie !");
        } catch (SQLException e) {
            JOptionPane.showMessageDialog(null,
                "Erreur de connexion : " + e.getMessage(),
                "Erreur", JOptionPane.ERROR_MESSAGE);
        }

        // === Configuration de la JFrame ===
        setTitle("Formulaire");
        setSize(650, 600);
        setDefaultCloseOperation(EXIT_ON_CLOSE);
        setLocationRelativeTo(null);
        setLayout(new BorderLayout());

        // === Panel titre ===
        titlePanel = new JPanel();
        titre = new JLabel("Formulaire d'inscription");
        titre.setFont(new Font("Serif", Font.BOLD, 28));
        titlePanel.add(titre);

        // === Panel formulaire (GridLayout 8x2) ===
        formPanel = new JPanel();
        formPanel.setBorder(new EmptyBorder(20, 40, 20, 40));
        formPanel.setLayout(new GridLayout(8, 2, 20, 20));

        lid         = new JLabel("IdMembre:");          id          = new JTextField();
        lprenom     = new JLabel("Prenom:");            prenom      = new JTextField();
        lnom        = new JLabel("Nom:");               nom         = new JTextField();
        ladresse    = new JLabel("Adresse:");           adresse     = new JTextField();
        ldateN      = new JLabel("Date de Naissance:"); dateN       = new JTextField();
        llieuN      = new JLabel("Lieu de Naissance:"); lieuN       = new JTextField();
        lsexe       = new JLabel("Sexe:");
        sexe        = new JComboBox<>();
        sexe.addItem("M"); sexe.addItem("F");
        lprofession = new JLabel("Profession:");        profession  = new JTextField();

        formPanel.add(lid);         formPanel.add(id);
        formPanel.add(lprenom);     formPanel.add(prenom);
        formPanel.add(lnom);        formPanel.add(nom);
        formPanel.add(ladresse);    formPanel.add(adresse);
        formPanel.add(ldateN);      formPanel.add(dateN);
        formPanel.add(llieuN);      formPanel.add(lieuN);
        formPanel.add(lsexe);       formPanel.add(sexe);
        formPanel.add(lprofession); formPanel.add(profession);

        // === Panel boutons ===
        buttonPanel = new JPanel();
        buttonPanel.setLayout(new FlowLayout(FlowLayout.CENTER, 40, 20));
        bAjouter  = new JButton("Ajouter");
        bModifier = new JButton("Modifier");
        bQuitter  = new JButton("Quitter");

        // ... (ActionListeners présentés dans le TP 11) ...

        buttonPanel.add(bAjouter);
        buttonPanel.add(bModifier);
        buttonPanel.add(bQuitter);

        add(titlePanel,  BorderLayout.NORTH);
        add(formPanel,   BorderLayout.CENTER);
        add(buttonPanel, BorderLayout.SOUTH);

        setVisible(true);
    }

    public static void main(String[] args) { new fenetrePrincipale(); }
}
\end{lstlisting}

\observation{Le bouton \textbf{Ajouter} exécute une requête
\texttt{INSERT INTO T\_Membre VALUES (?,?,?,?,?,?,?,?)} avec les valeurs
saisies dans les champs du formulaire. Le bouton \textbf{Modifier} exécute
une requête \texttt{UPDATE T\_Membre SET ... WHERE IdMembre=?}. Dans les
deux cas, \texttt{PreparedStatement} est utilisé pour éviter les injections
SQL et garantir la sécurité des données.}

\erreurFrequente{Oublier de fermer la connexion JDBC après chaque opération
entraîne des fuites de ressources qui peuvent épuiser le pool de connexions.
Utiliser \texttt{try-with-resources} avec \texttt{Connection},
\texttt{PreparedStatement} et \texttt{ResultSet} garantit leur fermeture
automatique.}

\probleme{Lors des premiers tests, le pilote JDBC n'était pas trouvé
(\texttt{ClassNotFoundException}). La solution a été d'ajouter le fichier
\texttt{.jar} du pilote (ex: \texttt{mysql-connector-java.jar}) au
\textit{classpath} du projet dans l'IDE. Une fois le pilote déclaré via
\texttt{Class.forName("com.mysql.cj.jdbc.Driver")}, la connexion a
fonctionné correctement.}

\conclusionTP{Ce TP marque l'aboutissement de la progression POO en
introduisant deux dimensions essentielles du développement Java professionnel :
les interfaces graphiques avec Swing et la persistance des données via JDBC.
La combinaison d'une \texttt{JFrame} bien structurée et d'une couche JDBC
sécurisée (\texttt{PreparedStatement}) constitue la base de nombreuses
applications de gestion desktop.}

% ════════════════════════════════════════════════════════════════════════════
%  TP 11 : Opérations CRUD avancées -- Swing & JDBC
% ════════════════════════════════════════════════════════════════════════════
\tpchapter{11}{Opérations CRUD Avancées -- Insertion et Recherche avec Swing \& JDBC}

\objectif{Approfondir l'intégration Swing/JDBC du TP~10 en implémentant deux
fonctionnalités clés : (1)~l'insertion des données du formulaire dans la table
\texttt{T\_Membre} via le bouton \textbf{Ajouter}, et (2)~la recherche et
l'affichage automatique d'un membre existant dans les champs du formulaire via
le bouton \textbf{Modifier}.}

% ── Question 1 ───────────────────────────────────────────────────────────────
\section{Question 1 -- Connexion et insertion (bouton \textbf{Ajouter})}

\enonce{Écrire le code Java qui, lors du clic sur le bouton \textbf{Ajouter},
établit une connexion avec la base de données, puis insère dans la table
\texttt{T\_Membre} les informations saisies par l'utilisateur dans le
formulaire Swing.}

\subsection{Logique d'implémentation}

Lors du clic sur le bouton \textbf{Ajouter}, le programme doit :

\begin{enumerate}[leftmargin=1.8cm, itemsep=3pt]
  \item Récupérer les valeurs saisies dans chaque composant Swing
        (\texttt{getText()} pour les \texttt{JTextField},
        \texttt{getSelectedItem()} pour le \texttt{JComboBox}).
  \item Établir la connexion JDBC via \texttt{DriverManager.getConnection()}.
  \item Construire un \texttt{PreparedStatement} avec la requête
        \texttt{INSERT INTO T\_Membre VALUES (?,?,?,?,?,?,?,?)}.
  \item Lier chaque paramètre \texttt{?} avec \texttt{setString()} ou
        \texttt{setInt()}.
  \item Exécuter l'insertion avec \texttt{executeUpdate()}.
  \item Afficher une confirmation à l'utilisateur et libérer les ressources.
\end{enumerate}

\subsection{Code Java -- Bouton \textbf{Ajouter}}

\begin{lstlisting}[caption={Gestionnaire du bouton Ajouter -- insertion JDBC}]
// Dans la methode actionPerformed de l'ActionListener du bouton "Ajouter"
if (e.getSource() == btnAjouter) {
    Connection conn = null;
    PreparedStatement pst = null;
    try {
        // 1. Recuperation des valeurs du formulaire
        String id        = txtIdMembre.getText().trim();
        String prenom    = txtPrenom.getText().trim();
        String nom       = txtNom.getText().trim();
        String adresse   = txtAdresse.getText().trim();
        String dateNais  = txtDateNaissance.getText().trim();
        String lieuNais  = txtLieuNaissance.getText().trim();
        String sexe      = cmbSexe.getSelectedItem().toString();
        String profession = txtProfession.getText().trim();

        // Validation basique : champs obligatoires
        if (id.isEmpty() || nom.isEmpty() || prenom.isEmpty()) {
            JOptionPane.showMessageDialog(this,
                "Les champs IdMembre, Nom et Prenom sont obligatoires.",
                "Saisie incomplete", JOptionPane.WARNING_MESSAGE);
            return;
        }

        // 2. Etablissement de la connexion JDBC
        Class.forName("com.mysql.cj.jdbc.Driver");
        conn = DriverManager.getConnection(
            "jdbc:mysql://localhost:3306/clubDB", "root", "");

        // 3. Preparation de la requete d'insertion
        String sql = "INSERT INTO T_Membre "
                   + "(IdMembre, Prenom, Nom, Adresse, DateNaissance, "
                   + " LieuNaissance, Sexe, Profession) "
                   + "VALUES (?, ?, ?, ?, ?, ?, ?, ?)";
        pst = conn.prepareStatement(sql);

        // 4. Liaison des parametres
        pst.setString(1, id);
        pst.setString(2, prenom);
        pst.setString(3, nom);
        pst.setString(4, adresse);
        pst.setString(5, dateNais);
        pst.setString(6, lieuNais);
        pst.setString(7, sexe);
        pst.setString(8, profession);

        // 5. Execution de l'insertion
        int rows = pst.executeUpdate();
        if (rows > 0) {
            JOptionPane.showMessageDialog(this,
                "Membre ajoute avec succes !",
                "Succes", JOptionPane.INFORMATION_MESSAGE);
            viderFormulaire(); // Reinitialiser les champs
        }
    } catch (SQLIntegrityConstraintViolationException ex) {
        JOptionPane.showMessageDialog(this,
            "Erreur : un membre avec cet IdMembre existe deja.",
            "Doublon", JOptionPane.ERROR_MESSAGE);
    } catch (ClassNotFoundException ex) {
        JOptionPane.showMessageDialog(this,
            "Pilote JDBC introuvable : " + ex.getMessage(),
            "Erreur pilote", JOptionPane.ERROR_MESSAGE);
    } catch (SQLException ex) {
        JOptionPane.showMessageDialog(this,
            "Erreur SQL : " + ex.getMessage(),
            "Erreur base de donnees", JOptionPane.ERROR_MESSAGE);
    } finally {
        // 6. Liberation des ressources
        try { if (pst  != null) pst.close();  } catch (SQLException ex) {}
        try { if (conn != null) conn.close(); } catch (SQLException ex) {}
    }
}
\end{lstlisting}

\screenshot{Screenshot 2026-05-TP11-ajouter.png}{Résultat -- insertion réussie et message de confirmation}

\observation{L'utilisation d'un \texttt{PreparedStatement} à la place d'un
\texttt{Statement} ordinaire est indispensable ici : elle protège contre les
injections SQL (un utilisateur malveillant pourrait sinon saisir du code SQL
dans un \texttt{JTextField}) et améliore les performances par précompilation
de la requête côté serveur.}

\note{La méthode \texttt{viderFormulaire()} (à implémenter) efface tous les
champs après une insertion réussie via \texttt{setText("")}, ce qui prépare le
formulaire pour une saisie suivante sans que l'utilisateur n'ait à effacer
manuellement chaque champ.}

\erreurFrequente{Oublier le bloc \texttt{finally} pour fermer
\texttt{Connection} et \texttt{PreparedStatement} provoque des fuites de
connexions qui peuvent saturer le nombre maximal de connexions autorisées par
le SGBD, rendant l'application inutilisable après quelques dizaines d'opérations.}

\probleme{Lors des premiers tests, une
\texttt{SQLIntegrityConstraintViolationException} était levée sans message
clair à l'utilisateur. En ajoutant un bloc \texttt{catch} spécifique pour ce
type d'exception et en affichant un \texttt{JOptionPane} explicatif, le
comportement est devenu convivial : l'utilisateur est informé qu'il tente
d'insérer un identifiant déjà existant.}

% ── Question 2 ───────────────────────────────────────────────────────────────
\section{Question 2 -- Recherche et affichage (bouton \textbf{Modifier})}

\enonce{Écrire le code Java pour le bouton \textbf{Modifier} qui permet de
rechercher un membre à partir de son identifiant (\texttt{IdMembre}) saisi par
l'utilisateur, puis d'afficher automatiquement les autres informations de ce
membre dans les zones de texte correspondantes.}

\subsection{Logique d'implémentation}

Lors du clic sur le bouton \textbf{Modifier}, le programme doit :

\begin{enumerate}[leftmargin=1.8cm, itemsep=3pt]
  \item Lire l'\texttt{IdMembre} saisi dans \texttt{txtIdMembre}.
  \item Exécuter une requête \texttt{SELECT * FROM T\_Membre WHERE IdMembre = ?}.
  \item Si un résultat est trouvé (\texttt{rs.next()} retourne \texttt{true}),
        remplir chaque composant Swing avec la valeur correspondante du
        \texttt{ResultSet}.
  \item Si aucun résultat n'est trouvé, afficher un message d'erreur via
        \texttt{JOptionPane}.
  \item Libérer toutes les ressources JDBC dans le bloc \texttt{finally}.
\end{enumerate}

\subsection{Code Java -- Bouton \textbf{Modifier}}

\begin{lstlisting}[caption={Gestionnaire du bouton Modifier -- recherche et remplissage JDBC}]
// Dans la methode actionPerformed de l'ActionListener du bouton "Modifier"
if (e.getSource() == btnModifier) {
    Connection conn = null;
    PreparedStatement pst = null;
    ResultSet rs = null;
    try {
        // 1. Lecture de l'identifiant saisi
        String id = txtIdMembre.getText().trim();
        if (id.isEmpty()) {
            JOptionPane.showMessageDialog(this,
                "Veuillez saisir un IdMembre pour effectuer la recherche.",
                "Champ manquant", JOptionPane.WARNING_MESSAGE);
            return;
        }

        // 2. Connexion JDBC
        Class.forName("com.mysql.cj.jdbc.Driver");
        conn = DriverManager.getConnection(
            "jdbc:mysql://localhost:3306/clubDB", "root", "");

        // 3. Requete SELECT avec filtre sur IdMembre
        String sql = "SELECT * FROM T_Membre WHERE IdMembre = ?";
        pst = conn.prepareStatement(sql);
        pst.setString(1, id);
        rs = pst.executeQuery();

        // 4. Remplissage du formulaire ou message d'erreur
        if (rs.next()) {
            txtPrenom.setText(rs.getString("Prenom"));
            txtNom.setText(rs.getString("Nom"));
            txtAdresse.setText(rs.getString("Adresse"));
            txtDateNaissance.setText(rs.getString("DateNaissance"));
            txtLieuNaissance.setText(rs.getString("LieuNaissance"));
            cmbSexe.setSelectedItem(rs.getString("Sexe"));
            txtProfession.setText(rs.getString("Profession"));

            JOptionPane.showMessageDialog(this,
                "Membre trouve. Vous pouvez modifier les informations.",
                "Membre charge", JOptionPane.INFORMATION_MESSAGE);
        } else {
            JOptionPane.showMessageDialog(this,
                "Aucun membre trouve avec l'identifiant : " + id,
                "Membre introuvable", JOptionPane.WARNING_MESSAGE);
            viderFormulaire();
        }
    } catch (ClassNotFoundException ex) {
        JOptionPane.showMessageDialog(this,
            "Pilote JDBC introuvable : " + ex.getMessage(),
            "Erreur pilote", JOptionPane.ERROR_MESSAGE);
    } catch (SQLException ex) {
        JOptionPane.showMessageDialog(this,
            "Erreur SQL : " + ex.getMessage(),
            "Erreur base de donnees", JOptionPane.ERROR_MESSAGE);
    } finally {
        // 5. Liberation des ressources (ordre inverse d'ouverture)
        try { if (rs   != null) rs.close();   } catch (SQLException ex) {}
        try { if (pst  != null) pst.close();  } catch (SQLException ex) {}
        try { if (conn != null) conn.close(); } catch (SQLException ex) {}
    }
}
\end{lstlisting}

\screenshot{Screenshot 2026-05-TP11-modifier.png}{Résultat -- champs du formulaire remplis automatiquement après recherche par IdMembre}

\observation{La séparation des deux actions en deux blocs \texttt{if}
distincts dans \texttt{actionPerformed} (un pour \texttt{Ajouter}, un pour
\texttt{Modifier}) est le patron standard des \texttt{ActionListener} Swing
multi-boutons. Cette approche est simple mais peut être refactorisée avec des
lambdas ou des \texttt{ActionListener} anonymes distincts pour chaque bouton.}

\note{Une fois les informations chargées depuis la base, l'utilisateur peut
modifier les champs du formulaire. Pour valider les modifications, il suffira
d'ajouter une troisième étape : un bouton ou une confirmation déclenchant une
requête \texttt{UPDATE T\_Membre SET ... WHERE IdMembre = ?} avec les nouvelles
valeurs -- ce qui constitue une évolution naturelle de ce TP.}

\erreurFrequente{Ne pas fermer le \texttt{ResultSet} avant de fermer le
\texttt{PreparedStatement} produit parfois des comportements indéfinis selon
le pilote JDBC. L'ordre correct de fermeture est toujours :
\texttt{ResultSet} $\to$ \texttt{PreparedStatement} $\to$ \texttt{Connection},
c'est-à-dire l'inverse de l'ordre d'ouverture.}

\probleme{Lors des tests, \texttt{cmbSexe.setSelectedItem()} ne sélectionnait
pas la valeur attendue si la chaîne stockée en base contenait une espace ou
était en minuscules. En appliquant \texttt{.trim().toUpperCase()} sur la valeur
lue depuis le \texttt{ResultSet} avant de l'affecter au \texttt{JComboBox},
la correspondance a fonctionné correctement dans tous les cas.}

\conclusionTP{Ce TP complète la boucle CRUD démarrée au TP~10 en ajoutant
une insertion robuste et une fonctionnalité de recherche dynamique. L'utilisation
systématique de \texttt{PreparedStatement}, de blocs \texttt{try/catch/finally}
et de messages \texttt{JOptionPane} produit une application fiable, sécurisée
et conviviale. Ces deux patterns -- insertion préparée et sélection par clé
primaire -- sont au cœur de toute application de gestion Java avec JDBC.}

% ────────────────────────────────────────────────────────────────────────────
\phantomsection
\addcontentsline{toc}{chapter}{Conclusion Générale}
\chapter*{Conclusion Générale}
\thispagestyle{fancy}

Ces dix travaux pratiques ont permis une progression structurée et complète
à travers les principaux piliers de la \textbf{Programmation Orientée Objet
en Java}. Chaque TP a enrichi les compétences acquises lors des sessions
précédentes, formant un parcours pédagogique cohérent allant des bases
syntaxiques jusqu'au développement d'applications graphiques connectées
à une base de données.

\vspace{0.6cm}

\begin{mdframed}[linecolor=bleuENSA, linewidth=1.5pt,
                 backgroundcolor=bleuClair,
                 innertopmargin=16pt, innerbottommargin=16pt,
                 innerleftmargin=18pt, innerrightmargin=18pt,
                 roundcorner=4pt,
                 topline=true, bottomline=true,
                 leftline=true, rightline=true]
\textbf{\color{bleuENSA}Récapitulatif des compétences acquises}

\vspace{8pt}
\renewcommand{\arraystretch}{1.35}
\begin{tabularx}{\linewidth}{@{}>{\bfseries}l X@{}}
\toprule
\rowcolor{bleuENSA!12}
\textbf{TPs} & \textbf{Compétences développées} \\
\midrule
TP 1--2  & Syntaxe Java, structures de contrôle, méthodes statiques, classes de base \\
TP 3--4  & Héritage, surcharge, classes abstraites, polymorphisme, liaison dynamique \\
TP 5     & Framework Collections, itérateurs, conversion, comparaison de performances \\
TP 6     & Hiérarchie d'exceptions, exceptions personnalisées, robustesse applicative \\
TP 7     & Programmation concurrente, synchronisation, communication inter-threads \\
TP 9     & Flux texte et binaire, sérialisation Java, persistance de données \\
TP 10    & Interfaces graphiques Swing, connexion JDBC, opérations CRUD \\
TP 11    & Insertion JDBC sécurisée, recherche par clé primaire, gestion des erreurs \\
\bottomrule
\end{tabularx}
\end{mdframed}

\vspace{0.8cm}

\conclusionTP{L'ensemble de ces travaux pratiques couvre les fondamentaux
incontournables du développement Java professionnel. La maîtrise de ces
concepts -- de l'encapsulation et de l'héritage jusqu'à la concurrence,
la persistance et les interfaces graphiques -- constitue le socle
indispensable pour aborder des frameworks modernes (Spring, Jakarta EE,
JavaFX) et des architectures logicielles avancées. Le langage Java, conçu
autour du principe «~écrire une fois, exécuter partout~», reste l'un des
langages les plus utilisés dans le monde professionnel.}

\vspace{0.5cm}

\begin{observationBox}
\textbf{\color{violetInfo}\faEye~~Perspectives d'approfondissement :}
Ces TPs ouvrent la voie à des thèmes avancés : les \textit{design patterns}
(Singleton, Factory, Observer), la programmation fonctionnelle avec les
\textit{lambda} et \textit{Streams} (Java 8+), le passage à JavaFX pour
des interfaces modernes, la gestion de la mémoire et du
\textit{garbage collection}, et les frameworks de test JUnit pour
garantir la qualité du code de façon automatisée.
\end{observationBox}

\vspace{0.8cm}
\begin{flushright}
  \begin{mdframed}[linecolor=grisBorder, linewidth=0.5pt,
    backgroundcolor=grisFond,
    innertopmargin=10pt, innerbottommargin=10pt,
    innerleftmargin=16pt, innerrightmargin=16pt,
    roundcorner=3pt]
    {\color{grisTexte}
      \textbf{\large LABSER Yahya}\\[3pt]
      \textit{IID1~~·~~ENSA Khouribga~~·~~2024/2025}
    }
  \end{mdframed}
\end{flushright}

\end{document}
